% THIS IS THE VERSION OF THE PAPER SUBMITTED TO EMSE
\documentclass[smallextended]{svjour3}

\newif\ifpaper\papertrue

% Paper-only packages
%\usepackage[pagewise]{lineno}
\usepackage{manyfoot}
\usepackage[most]{tcolorbox}
\usepackage[title]{appendix}
\usepackage{balance}
\usepackage{algorithm}
\usepackage{algorithmicx}
\usepackage[caption=false,font=normalsize,labelfont=sf,textfont=sf]{subfig}
\usepackage{stfloats}
\usepackage{graphicx}
\usepackage{xcolor}
\usepackage{tikz}
\usetikzlibrary{shapes.geometric, arrows, fit, shadows}
\usepackage[framemethod=tikz]{mdframed}

\setlength{\marginparwidth}{2cm}

% Shared header (packages, commands, conditionals)
% Shared header for llm-guidelines paper and website
% Requires: \newif\ifpaper (and \papertrue for paper) defined before \input

%%% Packages %%%

% Language and encoding
\usepackage[english]{babel}
\usepackage[utf8]{inputenc}
\usepackage[T1]{fontenc}

% URL/hyperlinks (url before hyperref, natbib before hyperref)
\usepackage[hyphens]{url}
\usepackage[square,numbers,sort&compress]{natbib}
\usepackage{hyperref}
\usepackage{csquotes}

% Formatting
\usepackage{xspace}
\usepackage[super]{nth}
\usepackage[inline]{enumitem}

% Math
\usepackage{amsmath}
\usepackage{amssymb}

% Tables
\usepackage{array}
\usepackage{tabularx}
\usepackage{booktabs}

% Code listings
\usepackage{listings}
\usepackage{textcomp} % for upquote=true

%%% Listings configuration %%%
\lstset{
  basicstyle=\ttfamily\small,
  showspaces=false,
  showstringspaces=false,
  showtabs=false,
  tabsize=2,
  breaklines=true,
  breakatwhitespace=true,
  aboveskip=0.5\baselineskip,
  belowskip=0.5\baselineskip,
  xleftmargin=0.5\parindent,
  xrightmargin=0pt,
  numbers=left,
  numbersep=8pt,
  upquote=true
}

\lstnewenvironment{plain}{\lstset{numbers=none}}{}

%%% Commands — identical in both %%%

\newcommand{\todo}[1]{{\textbf{TODO:}\ \textit{#1}}}
\newcommand*{\enq}[1]{\enquote{{\itshape#1}}}

%%% Commands — conditional %%%

% RFC 2119 keywords
\ifpaper
  \newcommand{\must}{\textsc{must}\xspace}
  \newcommand{\mustnot}{\textsc{must not}\xspace}
  \newcommand{\should}{\textsc{should}\xspace}
  \newcommand{\shouldnot}{\textsc{should not}\xspace}
\else
  \newcommand{\must}{\textbf{MUST}\xspace}
  \newcommand{\mustnot}{\textbf{MUST NOT}\xspace}
  \newcommand{\should}{\textbf{SHOULD}\xspace}
  \newcommand{\shouldnot}{\textbf{SHOULD NOT}\xspace}
\fi

% tl;dr label
\ifpaper
  \newcommand{\tldr}{\underline{\emph{tl;dr}}\xspace}
\else
  \newcommand{\tldr}{\textbf{\emph{tl;dr}}\xspace}
\fi

% Reporting location
\ifpaper
  \newcommand{\paper}{\emph{paper}\xspace}
  \newcommand{\supplementarymaterial}{\emph{supplementary material}\xspace}
\else
  \newcommand{\paper}{PAPER\xspace}
  \newcommand{\supplementarymaterial}{SUPPLEMENTARY MATERIAL\xspace}
\fi

% Icons for matrix/checklist
\ifpaper
  \newsavebox{\iconMbox}
  \savebox{\iconMbox}{\tikz[baseline=-0.5ex]\draw[fill=black] (0,0) circle (0.3em);}
  \newcommand{\iconM}{\usebox{\iconMbox}\xspace}
  \newsavebox{\iconSbox}
  \savebox{\iconSbox}{\tikz[baseline=-0.5ex]\draw[fill=gray!40, draw=gray!60] (0,0) circle (0.3em);}
  \newcommand{\iconS}{\usebox{\iconSbox}\xspace}
\else
  \newcommand{\iconM}{●\xspace}
  \newcommand{\iconS}{○\xspace}
\fi

% Framed environment for guideline summaries
\ifpaper
  \definecolor{framebg}{gray}{0.95}
  \newmdenv[
    leftmargin=0pt,
    rightmargin=0pt,
    usetwoside=false,
    innerleftmargin=6pt,
    innerrightmargin=6pt,
    innertopmargin=4pt,
    innerbottommargin=4pt,
    skipabove=0.5\baselineskip plus 1pt,
    skipbelow=0.5\baselineskip plus 2pt,
    linecolor=black,
    backgroundcolor=framebg,
    linewidth=2pt,
    roundcorner=0pt,
    nobreak=false,
    shadow=false,
    topline=false,
    bottomline=false,
    rightline=false,
    leftline=true,
    font=\itshape
  ]{framed}
\else
  \newenvironment{framed}{\begin{quote}\itshape}{\end{quote}}
\fi

% Cross-references — scope
\ifpaper
  \newcommand{\scope}{\hyperref[sec:scope]{\emph{Scope}}\xspace}
\else
  \newcommand{\scope}{\href{/scope/}{\emph{Scope}}\xspace}
\fi

% Cross-reference helpers with optional ID prefix
% Usage: \makepaperref{label}{Name}{ID} → "Name", \makepaperref[id]{label}{Name}{ID} → "ID: Name"
\newcommand{\makepaperref}[4][]{%
  \if\relax\detokenize{#1}\relax
    \hyperref[#2]{\emph{#3}}%
  \else
    \hyperref[#2]{\emph{#4: #3}}%
  \fi\xspace}
\newcommand{\makewebref}[4][]{%
  \if\relax\detokenize{#1}\relax
    \href{#2}{\emph{#3}}%
  \else
    \href{#2}{\emph{#4: #3}}%
  \fi\xspace}

% Cross-references — study types
\ifpaper
  \newcommand{\studytypes}{\hyperref[sec:study-types]{\emph{Study Types}}\xspace}
  \newcommand{\llmsforresearcher}{\hyperref[sec:llms-as-tools-for-software-engineering-researchers]{\emph{LLMs for Research}}\xspace}
  \newcommand{\annotators}[1][]{\makepaperref[#1]{sec:llms-as-annotators}{LLMs as Annotators}{S1}}
  \newcommand{\judges}[1][]{\makepaperref[#1]{sec:llms-as-judges}{LLMs as Judges}{S2}}
  \newcommand{\synthesis}[1][]{\makepaperref[#1]{sec:llms-for-synthesis}{LLMs for Synthesis}{S3}}
  \newcommand{\subjects}[1][]{\makepaperref[#1]{sec:llms-as-subjects}{LLMs as Subjects}{S4}}
  \newcommand{\llmsforengineers}{\hyperref[sec:llms-as-tools-for-software-engineers]{\emph{LLMs for SE}}\xspace}
  \newcommand{\llmusage}[1][]{\makepaperref[#1]{sec:studying-llm-usage-in-software-engineering}{Studying LLM Usage}{S5}}
  \newcommand{\newtools}[1][]{\makepaperref[#1]{sec:llms-for-new-software-engineering-tools}{LLMs for Tools}{S6}}
  \newcommand{\benchmarkingtasks}[1][]{\makepaperref[#1]{sec:benchmarking-llms-for-software-engineering-tasks}{Benchmarking LLMs}{S7}}
\else
  \newcommand{\studytypes}{\href{/study-types}{\emph{Study Types}}\xspace}
  \newcommand{\llmsforresearcher}{\href{/study-types/\#introduction-llms-as-tools-for-software-engineering-researchers}{\emph{LLMs for Research}}\xspace}
  \newcommand{\annotators}[1][]{\makewebref[#1]{/study-types\#llms-as-annotators}{LLMs as Annotators}{S1}}
  \newcommand{\judges}[1][]{\makewebref[#1]{/study-types\#llms-as-judges}{LLMs as Judges}{S2}}
  \newcommand{\synthesis}[1][]{\makewebref[#1]{/study-types\#llms-for-synthesis}{LLMs for Synthesis}{S3}}
  \newcommand{\subjects}[1][]{\makewebref[#1]{/study-types\#llms-as-subjects}{LLMs as Subjects}{S4}}
  \newcommand{\llmsforengineers}{\href{/study-types/\#introduction-llms-as-tools-for-software-engineers}{\emph{LLMs for SE}}\xspace}
  \newcommand{\llmusage}[1][]{\makewebref[#1]{/study-types\#studying-llm-usage-in-software-engineering}{Studying LLM Usage}{S5}}
  \newcommand{\newtools}[1][]{\makewebref[#1]{/study-types\#llms-for-new-software-engineering-tools}{LLMs for Tools}{S6}}
  \newcommand{\benchmarkingtasks}[1][]{\makewebref[#1]{/study-types\#benchmarking-llms-for-software-engineering-tasks}{Benchmarking LLMs}{S7}}
\fi

% Cross-references — guidelines
\ifpaper
  \newcommand{\guidelines}{\hyperref[sec:guidelines]{\emph{Guidelines}}\xspace}
  \newcommand{\usagerole}[1][]{\makepaperref[#1]{sec:declare-llm-usage-and-role}{Usage and Role}{G1}}
  \newcommand{\modelversion}[1][]{\makepaperref[#1]{sec:report-model-version-configuration-and-customizations}{Version and Configuration}{G2}}
  \newcommand{\toolarchitecture}[1][]{\makepaperref[#1]{sec:report-tool-architecture-beyond-models}{Architecture}{G3}}
  \newcommand{\humanvalidation}[1][]{\makepaperref[#1]{sec:use-human-validation-for-llm-outputs}{Human Validation}{G5}}
  \newcommand{\prompts}[1][]{\makepaperref[#1]{sec:report-prompts-their-development-and-interaction-logs}{Prompts and Logs}{G4}}
  \newcommand{\openllm}[1][]{\makepaperref[#1]{sec:use-an-open-llm-as-a-baseline}{Open LLMs}{G6}}
  \newcommand{\benchmarksmetrics}[1][]{\makepaperref[#1]{sec:use-suitable-baselines-benchmarks-and-metrics}{Benchmarks and Metrics}{G7}}
  \newcommand{\limitationsmitigations}[1][]{\makepaperref[#1]{sec:report-limitations-and-mitigations}{Limitations and Mitigations}{G8}}
\else
  \newcommand{\guidelines}{\href{/guidelines}{\emph{Guidelines}}\xspace}
  \newcommand{\usagerole}[1][]{\makewebref[#1]{/guidelines\#declare-llm-usage-and-role}{Usage and Role}{G1}}
  \newcommand{\modelversion}[1][]{\makewebref[#1]{/guidelines\#report-model-version-configuration-and-customizations}{Version and Configuration}{G2}}
  \newcommand{\toolarchitecture}[1][]{\makewebref[#1]{/guidelines\#report-tool-architecture-beyond-models}{Architecture}{G3}}
  \newcommand{\humanvalidation}[1][]{\makewebref[#1]{/guidelines\#use-human-validation-for-llm-outputs}{Human Validation}{G5}}
  \newcommand{\prompts}[1][]{\makewebref[#1]{/guidelines\#report-prompts-their-development-and-interaction-logs}{Prompts and Logs}{G4}}
  \newcommand{\openllm}[1][]{\makewebref[#1]{/guidelines\#use-an-open-llm-as-a-baseline}{Open LLMs}{G6}}
  \newcommand{\benchmarksmetrics}[1][]{\makewebref[#1]{/guidelines\#use-suitable-baselines-benchmarks-and-metrics}{Benchmarks and Metrics}{G7}}
  \newcommand{\limitationsmitigations}[1][]{\makewebref[#1]{/guidelines\#report-limitations-and-mitigations}{Limitations and Mitigations}{G8}}
\fi

% Section formatting
\ifpaper
  \newcommand{\guidelinesubsubsection}[1]{\subsubsection{\bfseries #1:}}
  \newcommand{\studytypesubsection}[1]{\subsubsection{#1}}
  \newcommand{\studytypeparagraph}[1]{\paragraph{\bfseries #1:}}
  \newcommand{\scopeparagraph}[1]{\paragraph{\bfseries #1:}}
\else
  \newcommand{\guidelinesubsubsection}[1]{\subsubsection{#1}}
  \newcommand{\studytypesubsection}[1]{\subsection{#1}}
  \newcommand{\studytypeparagraph}[1]{\subsubsection{#1}}
  \newcommand{\scopeparagraph}[1]{\subsubsection{#1}}
\fi
 
% Paper-only: lstset color overrides (requires xcolor loaded above)
\lstset{
  keywordstyle=\color{blue},
  identifierstyle=\color{black},
  commentstyle=\color{gray},
  stringstyle=\color{gray},
  emphstyle=\color{red},
}

% Paper-only settings
\renewcommand{\sectionautorefname}{Section}
\renewcommand{\subsectionautorefname}{Section}
\renewcommand{\subsubsectionautorefname}{Section}

% prevent single lines at the start of a paragraph (Schusterjungen)
\clubpenalty = 10000
% prevent single lines at the end of a paragraph (Hurenkinder)
\widowpenalty = 10000 \displaywidowpenalty = 10000

% Set title
\makeatletter
\def\makeheadbox{%
  \hbox to0pt{%
    \vbox{\baselineskip=10dd
      \hrule
      \hbox to\hsize{%
        \vrule\kern3pt
        \vbox{\kern3pt
          % ---- customise here -----------------
          \hbox{\bfseries Empirical Software Engineering}%
          \hbox{(submitted manuscript)}%
          % --------------------------------------
          \kern3pt}%
        \hfil\kern3pt\vrule}%
      \hrule}%
    \hss}}
\makeatother

\date{Submitted: 26/08/2025\\Revised: 27/02/2026}

\title{Guidelines for Empirical Studies in Software~Engineering involving Large~Language~Models}

\author{Sebastian~Baltes \and Florian~Angermeir \and Chetan~Arora \and Marvin~Muñoz~Barón \and Chunyang~Chen \and Lukas~Böhme \and Fabio~Calefato \and Neil~Ernst \and Davide~Falessi \and Brian~Fitzgerald \and Davide~Fucci \and Junda~He \and Christoph~Treude \and Marcos~Kalinowski \and Stefano~Lambiase \and Daniel~Russo \and Mircea~Lungu \and Cristina~Martinez~Montes \and Lutz~Prechelt \and Paul~Ralph  \and Rijnard~van~Tonder  \and Stefan~Wagner}

\institute{%
  Sebastian~Baltes \at
  Heidelberg University, Germany\\
  \email{sebastian.baltes@uni-heidelberg.de}\\
  (corresponding author)
  \and
  Florian~Angermeir \at
  fortiss, Germany\\
  Blekinge Institute of Technology, Sweden\\
  \email{angermeir@fortiss.org}
  \and
  Chetan~Arora \at
  Monash University, Australia\\
  \email{chetan.arora@monash.edu}
  \and
  Marvin~Muñoz~Barón  \and Chunyang~Chen \at
  Technical University of Munich, Germany\\
  \email{\{marvin.munoz-baron, chun-yang.chen\}@tum.de}
  \and
  Lukas~Böhme \at
  Hasso-Plattner-Institut, Germany\\
  University of Potsdam, Germany\\
  \email{lukas.boehme@hpi.de}
  \and
  Fabio~Calefato \at
  University of Bari, Italy\\
  \email{fabio.calefato@uniba.it}
  \and
  Neil~Ernst \at
  University of Victoria, Canada\\
  \email{nernst@uvic.ca}
  \and
  Davide~Falessi \at
  University of Rome Tor Vergata, Italy\\
  \email{falessi@ing.uniroma2.it}
  \and
  Brian~Fitzgerald \at
  Lero, Ireland\\
  University of Limerick, Ireland\\
  \email{brian.fitzgerald@ul.ie}
  \and
  Davide~Fucci \at
  Blekinge Institute of Technology, Sweden\\
  \email{davide.fucci@bth.se}
  \and
  Junda~He \and Christoph~Treude \at
  Singapore Management University, Singapore\\
  \email{\{jundahe, ctreude\}@smu.edu.sg}
  \and
  Marcos~Kalinowski \at
  PUC Rio de Janeiro, Brazil\\
  \email{kalinowski@inf.puc-rio.br}
  \and
  Stefano~Lambiase \and Daniel~Russo \at
  Aalborg University in Copenhagen, Denmark\\
  \email{\{stla, daniel.russo\}@cs.aau.dk}
  \and
  Mircea~Lungu \at
  IT University of Copenhagen, Denmark\\
  \email{mlun@itu.dk}
  \and
  Cristina~Martinez~Montes \at
  Chalmers University of Technology, Sweden\\
  University of Gothenburg, Sweden\\
  \email{montesc@chalmers.se}
  \and
  Lutz~Prechelt \at
  Freie Universität Berlin, Germany\\
  \email{prechelt@inf.fu-berlin.de}
  \and
  Paul~Ralph \at
  Dalhousie University, Canada\\
  \email{paulralph@dal.ca}
  \and
  Rijnard~van~Tonder \at
  Independent Researcher, Antigua and Barbuda\\
  \email{rvantonder@gmail.com}
  \and
  Stefan~Wagner \at
  Technical University of Munich, Germany\\
  \email{stefan.wagner@tum.de}
}

\authorrunning{Baltes et al.}
\titlerunning{Guidelines for Empirical Studies in SE involving LLMs}

\begin{document}
%\linenumbers

\maketitle

\begin{abstract}
Large Language Models (LLMs) are now ubiquitous in software engineering (SE)
research and practice, yet their non-determinism, opaque training data, and
rapidly evolving models threaten the reproducibility and replicability of
empirical studies.
We address this challenge through a collaborative effort of 22 researchers,
presenting a taxonomy of seven study types that organizes the landscape of
LLM involvement in SE research, together with eight guidelines for designing
and reporting such studies.
Each guideline distinguishes requirements (\must) from recommended practices
(\should) and is contextualized by the study types it applies to.
Our guidelines recommend that researchers:
(1)~declare LLM usage and role;
(2)~report model versions, configurations, and customizations;
(3)~document the tool architecture beyond the model;
(4)~disclose prompts, their development, and interaction logs;
(5)~validate LLM outputs with humans;
(6)~include an open LLM as a baseline;
(7)~use suitable baselines, benchmarks, and metrics; and
(8)~articulate limitations and mitigations.
We complement the guidelines with an applicability matrix mapping guidelines
to study types and a reporting checklist for authors and reviewers.
We maintain the study types and guidelines online as a living resource for the
community to use and shape
(\href{https://llm-guidelines.org/}{llm-guidelines.org}).
 \end{abstract}

\keywords{software engineering, large language models, empirical research, meta-science, guidelines}

\section{Introduction}
\label{sec:introduction}


Since the release of ChatGPT in November 2022, large language models (LLMs) have rapidly become a major focus of software engineering (SE) research~\cite{10.1145/3695988}, yet the reproducibility and replicability of empirical studies involving LLMs remains uncertain.
Recent findings indicate a low reproducibility of SE studies involving LLMs~\cite{DBLP:journals/corr/abs-2510-25506}.
Three characteristics of LLMs make this particularly challenging:
their inherent non-determinism causes variability across runs, with slight changes leading to significantly different results~\cite{DBLP:journals/corr/abs-2412-12509, DBLP:conf/naacl/SongWLL25, DBLP:conf/nips/AgarwalSCCB21, bjarnason2026randomnessagenticevals};
commercial models evolve beyond version identifiers, so reported performance can change over time~\cite{DBLP:journals/corr/abs-2307-09009};
and even for ``open'' models, training data and fine-tuning details often remain undisclosed~\cite{Gibney2024}.
Moreover, prompt formatting choices alone can shift accuracy by tens of percentage points~\cite{DBLP:conf/iclr/Sclar0TS24}, and configured parameters such as temperature affect output variability~\cite{renze2024temperature}.
Hence, not reporting these settings directly affects reproducibility.

Although the SE research community has developed guidelines for conducting and reporting specific types of empirical studies such as controlled experiments (e.g.,~\cite{DBLP:books/sp/WohlinRHORW24, DBLP:books/sp/08/SSS2008}), their replications (e.g.,~\cite{DBLP:journals/tse/SantosVOJ21}), or empirical studies in general (e.g., the \emph{ACM SIGSOFT Empirical Standards}~\cite{ralph2021empiricalstandardssoftwareengineering}), none of these address the LLM-specific aspects described above.
Previously, a position paper highlighted these issues~\cite{DBLP:conf/wsese/0001BFB25}, but there was no comprehensive community-developed guidance for designing and reporting empirical studies involving LLMs in SE.

Therefore, we present community-developed guidelines for conducting and reporting studies involving LLMs in SE research, co-developed by 22 researchers.
After outlining our \scope, we introduce a taxonomy of \studytypes, then present eight \guidelines.
We complement these with an applicability matrix mapping guidelines to study types and a reporting checklist for authors and reviewers.
For each study type and guideline, we identify relevant examples, both within and outside of SE research.
We maintain the guidelines online as a living resource for the community to use and shape.\footnote{\url{https://llm-guidelines.org}}
 
\section{Scope of Guidelines}
\label{sec:scope}

\scopeparagraph{Software Engineering as our Target Discipline}

We target SE research because existing cross-disciplinary guidelines do not address its needs. For instance, Gallifant et al.'s~\cite{Gallifant2025} guidelines for LLM use in healthcare research include discipline-specific items irrelevant to most SE studies (e.g., clinical-care usability) while lacking guidance on tool architectures, benchmarking, and code-specific evaluation that SE research requires.
Moreover, SE research employs a wide variety of empirical methods~\cite{ralph2021empiricalstandardssoftwareengineering, DBLP:books/sp/WohlinRHORW24}. Our work builds on \citeauthor{sallou2024breaking}'s vision paper on mitigating validity threats in LLM-based SE research~\cite{sallou2024breaking} and our own position paper~\cite{DBLP:conf/wsese/0001BFB25}, providing more detailed advice for a wider variety of research methods organized into a taxonomy of \studytypes (Section~\ref{sec:study-types}).

\scopeparagraph{Focus on Text-Based Use Cases}

While multi-modal foundation models that use or generate images, audio, or video may also support SE research and practice, we focus on textual use cases of LLMs (e.g., in natural language or programming languages). Many of our guidelines---particularly those concerning model reporting, prompt documentation, and reproducibility---are likely applicable to multi-modal settings as well, but we leave their explicit validation to future work.

\scopeparagraph{Focus on Direct Development or Research Support}

While researchers may use LLMs for many peripheral tasks (e.g., proof-reading, spell-checking, translation), our guidelines focus on their direct role in empirical research and engineering practice.
For engineers, we focus on the use of LLMs to automate SE tasks, that is, artificial intelligence (AI) for software engineering (AI4SE) (see~\llmsforengineers).
This includes agentic systems that autonomously plan and execute multi-step tasks using LLMs (see~\toolarchitecture).
For researchers, we focus on the use of LLMs to automate empirical research tasks such as data collection, processing, or analysis (see~\llmsforresearcher).

\scopeparagraph{Researchers as our Target Audience}

Our guidelines are intended to help SE researchers design, plan, conduct,
and report empirical studies involving LLMs, and to support scholarly peer
review of such studies. Each guideline includes an \emph{Advice for
Reviewers} subsection with targeted assessment suggestions.
Our guidelines focus on what to report and how; they complement but do not replace methodological guidance for designing specific types of empirical studies.

\scopeparagraph{How to Navigate this Paper}

This paper is structured to support different reading strategies depending on the reader's goal.
\emph{Researchers planning a new study} may start with the taxonomy of
study types (see \studytypes) to identify which types apply
to their planned work, then consult Table~\ref{tab:guideline-matrix} to
determine which guidelines are requirements (\must) and which are
recommendations (\should) for those study types. Each guideline section
opens with a \tldr summary in a shaded box, allowing
readers to quickly assess relevance before reading the full text.
\emph{Researchers writing up results} may prefer to start with the
checklist in Appendix~\ref{sec:checklist}, which organizes actionable items
by typical paper sections (Introduction, Research Design and Methods, Results, etc.).
Table~\ref{tab:guidelines} provides a quick reference to all eight
guidelines, and Table~\ref{tab:rationale-recommendations} maps each
guideline's rationale to its key recommendations.
\emph{Reviewers} can use the \emph{Advice for Reviewers} subsection at the
end of each guideline for targeted guidance on assessing manuscripts.
Table~\ref{tab:guideline-matrix} helps reviewers identify which guidelines
apply to the study type under review. 
\section{Methodology}
\label{sec:methodology}

This project was initiated at the 2024 meeting of the \emph{International Software Engineering Research Network} (ISERN) in Barcelona, Spain, where the first and last author organized a session on guidelines for empirical studies involving LLMs.
These discussions led to a position paper at \emph{WSESE 2025}~\citep{DBLP:conf/wsese/0001BFB25}, which was presented at the \emph{\nth{2} Copenhagen Symposium on Human-Centered Software Engineering AI} (CHASEAI 2024), forming a workstream to collaboratively refine and extend the preliminary guidelines.
Our process follows an established tradition of developing reporting guidelines through expert collaboration, as exemplified by CONSORT for clinical trials~\citep{Begg1996} and, more recently, TRIPOD-LLM for LLM-based prediction studies in medicine~\citep{Gallifant2025}.

We held bi-weekly meetings, in which we decided on the structure of the study type taxonomy and the guidelines.
During these discussions, we added a new study type (\synthesis), refined and extended the guideline sections (in particular \modelversion, \toolarchitecture, \prompts) and covered \benchmarksmetrics as well as \limitationsmitigations.

Then, at least two co-authors were assigned to each study type and guideline to refine and extend the content, followed by a peer review and refinement phase.
Afterwards, the first author reviewed all sections and discussed potential changes with the topic leads.
A requirement for co-authorship was that each author contributed or reviewed content and, most importantly, read the complete guidelines and confirmed that they stand behind all recommendations, not only their own contributions.

To select illustrative examples for each study type and guideline, co-author pairs independently searched for relevant papers using academic databases.
This initial pool was extended based on suggestions from other co-authors.
Identified papers were assessed by one author and subsequently reviewed by a second author; papers that did not add new insights beyond those already covered were not included.
Multiple review rounds were then carried out on the guidelines as a whole, during which selected papers were cross-checked by additional authors and, where necessary, further examples were added or replaced.
The examples are intended to be illustrative, not the result of a systematic review; the involvement of multiple authors at different stages helped mitigate individual selection bias.

After completing the internal review, we invited external experts in empirical SE to review the study types and guidelines.
During the review process of the \emph{Empirical Software Engineering} journal, the structure and content were improved based on reviewers' suggestions.
We self-assessed our guidelines against the quality dimensions of the \emph{SIGSOFT Empirical Standards}~\citep{ralph2021empiricalstandardssoftwareengineering}; this assessment is reported in the conclusion.


\section{Study Types}
\label{sec:study-types}

Empirical guidelines are needed for SE studies involving LLMs to ensure the validity and reproducibility of their results.
However, such guidelines must be tailored to different study types that each pose unique challenges.
%Therefore, understanding the classification of these studies is essential for developing appropriate guidelines.
Therefore, we created a taxonomy of study types to contextualize our recommendations.
%Note that our \guidelines refer to the \studytypes but not the other way around.
We use the term \emph{study type} to refer to categories of LLM involvement in empirical research, rather than to research methodologies (e.g., experiments, case studies). A single empirical study may involve multiple such study types.
Each study type section starts with a \textbf{description}, followed by \textbf{examples} from the SE research community and beyond. The \textbf{advantages} and \textbf{challenges} of using LLMs are discussed in a summary subsection at the end of this section, synthesizing cross-cutting themes across all study types.
\ifpaper
See Table~\ref{tab:study-types} for an overview of study types.
\fi
 
%% This sub-table of contents is organizational overkill. 
%% I think it's  a good additional since the guidlines are already a wall of text.

\begin{table}[tb]
\caption{Overview of study types.}
\label{tab:study-types}
\begin{tabular}{@{}lll@{}}
\toprule
\textbf{ID} & \textbf{Section} & \textbf{Topic} \\
\midrule
     & 4.1   & \textbf{\hyperref[sec:llms-as-tools-for-software-engineering-researchers]{LLMs as Tools for Software Engineering Researchers}} \\
S1   & 4.1.1 & \quad \hyperref[sec:llms-as-annotators]{LLMs as Annotators} \\
S2   & 4.1.2 & \quad \hyperref[sec:llms-as-judges]{LLMs as Judges} \\
S3   & 4.1.3 & \quad \hyperref[sec:llms-for-synthesis]{LLMs for Synthesis} \\
S4   & 4.1.4 & \quad \hyperref[sec:llms-as-subjects]{LLMs as Subjects} \\
\midrule
     & 4.2   & \textbf{\hyperref[sec:llms-as-tools-for-software-engineers]{LLMs as Tools for Software Engineers}} \\
S5   & 4.2.1 & \quad \hyperref[sec:studying-llm-usage-in-software-engineering]{Studying LLM Usage in Software Engineering} \\
S6   & 4.2.2 & \quad \hyperref[sec:llms-for-new-software-engineering-tools]{LLMs for New Software Engineering Tools} \\
S7   & 4.2.3 & \quad \hyperref[sec:benchmarking-llms-for-software-engineering-tasks]{Benchmarking LLMs for Software Engineering Tasks} \\
\bottomrule
\end{tabular}
\end{table}

\subsection{LLMs as Tools for Software Engineering Researchers}
\label{sec:llms-as-tools-for-software-engineering-researchers}

LLMs can serve as tools to help researchers conduct empirical studies.
They can automate various tasks such as data collection, pre-processing, and analysis.
For example, LLMs can apply predefined coding guides to qualitative datasets (\annotators), assess the quality of software artifacts (\judges), generate summaries of research papers (\synthesis), or simulate human behavior in empirical studies (\subjects). If LLMs could complete these tasks well enough, they would reduce the time and effort required to conduct a study. However, these applications each pose challenges to validity and reproducibility.
 
\studytypesubsection{LLMs as Annotators (S1)}
\label{sec:llms-as-annotators}

\studytypeparagraph{Description}
In qualitative data analysis, manually annotating (``coding'') natural language text (e.g., requirements, interview transcripts, open-ended survey responses) is a time-consuming manual process~\cite{DBLP:journals/ase/BanoHZT24}. LLMs can augment human coding, suggest new codes, and label artifacts based on a pre-defined coding guide much faster than humans can~\cite{DBLP:conf/chi/HeHDRH24}. 
%(see \synthesis), or even automate the entire qualitative data analysis process.
The extent to which, or under what conditions, LLMs can perform these tasks \textit{effectively} remains an open research question~\cite{DBLP:conf/msr/AhmedDTP25}. Indeed, measuring their effectiveness is practically and philosophically challenging. From an interpretivist philosophical perspective, one cannot measure the quality of analysis by comparing one (human or machine) analyst's work to another. From a realist perspective, triangulating across multiple human judges, LLM judges, and other data sources (e.g., whether a pull request is marked as having resolved an issue) improves confidence in the findings but does not prove that any one judge is valid.

\studytypeparagraph{Examples}
\citeauthor{Huang2023Enhancing}~\cite{Huang2023Enhancing} utilized multiple LLMs for joint annotation of mobile application reviews. 
They used three models of comparable size with an absolute majority voting rule (i.e., a label is only accepted if it receives more than half of the total votes from the models). This approach slightly outperformed the best individual model tested. 
Meanwhile, \citeauthor{DBLP:conf/msr/AhmedDTP25}~\cite{DBLP:conf/msr/AhmedDTP25} examined LLMs as annotators in SE research across five datasets, six LLMs, and ten annotation tasks.
They found that model-model agreement strongly correlates with human-model agreement; models performed poorly in tasks where humans also frequently disagreed. They proposed to use model confidence scores to identify specific samples that could be safely delegated to LLMs, potentially reducing human annotation effort without compromising inter-rater agreement.
%% Reliabiltiy is a necessary condition for, but does not constitute, validity. We cannot infer that a statement is true just because several LLMs agree on it.  
%Conversely, if multiple LLMs reach similar solutions independently, then LLMs are likely suitable for the annotation task.
 
\studytypesubsection{LLMs as Judges (S2)}
\label{sec:llms-as-judges}

\studytypeparagraph{Description}

LLMs can rate properties of software artifacts (e.g. assess code readability, adherence to coding standards, or comment quality) or sort multiple solutions by some attribute. These kinds of judgments are distinct from annotating unstructured text (see \annotators).

\studytypeparagraph{Example(s)}

\citeauthor{DBLP:conf/re/LubosFTGMEL24}~\cite{DBLP:conf/re/LubosFTGMEL24} leveraged \href{https://www.llama.com/llama2/}{Llama-2} to evaluate the quality of software requirements statements. 
They prompted the LLM with the text below, where the words in braces reflect the study parameters:

\begin{plain}
Your task is to evaluate the quality of a software requirement.
Evaluate whether the following requirement is {quality_characteristic}.
{quality_characteristic} means: {quality_characteristic_explanation}
The evaluation result must be: 'yes' or 'no'.
Request: Based on the following description of the project: {project_description}
Evaluate the quality of the following requirement: {requirement}.
Explain your decision and suggest an improved version.
\end{plain}

They evaluated LLM output against expert human judges and found moderate agreement for simple requirements and poor agreement for more complex requirements. In contrast, \citeauthor{wang2025multimodalrequirementsdatabasedacceptance}~\cite{wang2025multimodalrequirementsdatabasedacceptance} used an LLM to generate acceptance criteria for user stories, and provided a rubric to an LLM to judge the generated acceptance criteria on interpretable scales (0 to 4). 
 
\studytypesubsection{LLMs for Synthesis (S3)}
\label{sec:llms-for-synthesis}

\studytypeparagraph{Description}

Unlike annotation (see \annotators), which focuses on categorizing or labeling individual data points, synthesis refers to the process of integrating and interpreting information from multiple sources to generate higher-level insights, identify patterns across datasets, and develop conceptual frameworks or theories. LLMs may be able to support synthesis tasks in SE research by processing and distilling information from qualitative data sources. 
Although synthesis in the preceding notion refers to abstraction and interpretation across multiple data sources, the term is sometimes also used to refer to generating synthetic content (e.g., source code, bug-fix pairs, requirements, etc.) that are then used in downstream tasks to train, fine-tune, or evaluate existing models or tools.
In this case, the synthesis is done primarily using the LLM and its training data; the input is limited to basic instructions and examples.

\studytypeparagraph{Example(s)} 

Published examples of applying LLMs for synthesis in SE remain scarce; however, some recent work in other domains is instructive~\cite{DBLP:journals/ase/BanoHZT24}.
\citeauthor{barros2024largelanguagemodelqualitative}~\cite{barros2024largelanguagemodelqualitative} conducted a systematic mapping study on using LLMs for qualitative research and found that most studies focused on the healthcare and social sciences. In these studies, LLMs supported qualitative methods, such as grounded theory and thematic analysis, by aiding in pattern identification.
%They identified examples in domains such as healthcare and social sciences (see, e.g., \cite{de2024performing,mathis2024inductive}) in which LLMs were used to support qualitative analysis, including grounded theory and thematic analysis.
%Overall, the findings highlight the successful generation of preliminary coding schemes from interview transcripts, later refined by human researchers, along with support for pattern identification.This approach was reported not only to expedite the initial coding process but also to allow researchers to focus more on higher-level analysis and interpretation.
%This particular paper suggests using LLMs to support tasks such as initial coding and theme identification while conservatively reserving interpretative or creative processes for human analysts.
In SE, de Morais Leca et al. \cite{leça2024applicationsimplicationslargelanguage} explored how LLMs have been applied for qualitative data analysis (QDA) and proposed general strategies and guidelines for their application. Ornelas et al. \cite{ornelas2025llm} complemented this perspective by studying the opportunities and limitations of introducing LLM-based support into QDA, and by formulating recommendations for embedding human–AI collaboration across the thematic analysis phases. Building on these, subsequent work has proposed hybrid frameworks combining LLM support with human-led QDA. Rashid et al. \cite{DBLP:journals/corr/abs-2402-01386} designed an LLM-driven multi-agent system that integrates AI with human decision-making to automate qualitative data analysis methods. Their system generated initial codes, developed themes, and summarized text. Similarly, Montes et al. \cite{DBLP:journals/corr/abs-2510-18456} compared the performance of humans and LLMs in coding, theme development, definition, and refinement, creating guidelines for a hybrid-LLM framework.
Finally, \citeauthor{DBLP:journals/corr/abs-2506-21138}'s work is an example of using LLMs to create synthetic datasets.
They present an approach to generate synthetic requirements, showing that they \enq{can match or surpass human-authored requirements for specific classification tasks}~\cite{DBLP:journals/corr/abs-2506-21138}.
 
% Reverted back to "subject", as participant seems to be primarily used for humans: see  https://dictionary.apa.org/subject and https://dictionary.apa.org/participant
\studytypesubsection{LLMs as Subjects (S4)}
\label{sec:llms-as-subjects}

\studytypeparagraph{Description}

In empirical studies, data is collected from participants through methods such as surveys, interviews, or controlled experiments.
LLMs can serve as virtual \emph{subjects} by simulating human behavior and interactions. If LLMs can generate responses that approximate those of human participants, they could be valuable for research involving user interactions, collaborative coding environments, and software usability assessments~\cite{ZHAO2025101167}.
To achieve this, prompt engineering techniques are widely employed; for instance, the \textit{Personas Pattern}~\cite{DBLP:journals/corr/abs-2308-07702} involves tailoring LLM responses to align with predefined profiles or roles that emulate specific user archetypes.
To serve as virtual subjects, generated responses should be indistinguishable from human-produced texts, consistent with the attitudes and sociodemographic information of the conditioning context (e.g., junior vs.\ senior developers), naturally aligned with the form, tone, and content of the simulated scenario, and reflect patterns in relationships between ideas, demographics, and behavior observed in comparable human data~\cite{DBLP:journals/corr/abs-2209-06899}.

\studytypeparagraph{Example(s)}

\citeauthor{DBLP:journals/ipm/XuSRGPLSH24}~\cite{DBLP:journals/ipm/XuSRGPLSH24} compiled a list of ways LLMs can support social science research, some of which transfer to empirical SE research. For example, LLMs can emulate human responses and behaviors in simulated interviews and focus groups~\citeauthor{DBLP:journals/ase/GerosaTSS24}~\cite{DBLP:journals/ase/GerosaTSS24}. Similarly, \citeauthor{bano2025doessoftwareengineerlook}~\cite{bano2025doessoftwareengineerlook}, investigated biases in LLM-generated candidate profiles in SE recruitment processes. They found biases favoring male candidates, lighter skin tones, and slim physiques, particularly for senior roles. LLMs may be able to simulate end-user feedback and behavior in usability studies, identify usability issues and offering suggestions for improvement based on predefined user personas.
 
\subsection{LLMs as Tools for Software Engineers}
\label{sec:llms-as-tools-for-software-engineers}

LLM-based assistants have become a widely adopted tool for software engineers~\cite{stackoverflow2025survey}, supporting them in various tasks such as code generation and debugging.
Researchers have studied how software engineers use LLMs (\llmusage), developed new tools that integrate LLMs (\newtools), and benchmarked LLMs for software engineering tasks (\benchmarkingtasks).
%Like in the previous section, we outline the advantages that studying LLMs in these contexts brings, but also point to specific challenges.
 
\studytypesubsection{Studying LLM Usage in Software Engineering (S5)}
\label{sec:studying-llm-usage-in-software-engineering}

\studytypeparagraph{Description}

Studying how software engineers use LLMs and LLM-based tools is important for understanding the current state of practice in SE.
Researchers can observe software engineers' usage of LLM-based tools in the field, or study if and how they adopt such tools, their usage patterns, as well as perceived benefits and challenges.
Surveys, interviews, observational studies, or analysis of usage logs can provide insights into how LLMs are integrated into development processes, how they influence decision making, and what factors affect their acceptance and effectiveness. 
Such studies can inform improvements for existing LLM-based tools, motivate the design of novel tools, or derive best practices for LLM-assisted software engineering.
They can also uncover risks or deficiencies of existing tools.

\studytypeparagraph{Example(s)}

Based on a convergent mixed-methods study, \citeauthor{russo2024navigating} has found that early adoption of generative AI by software engineers is primarily driven by compatibility with existing workflows~\cite{russo2024navigating}.
\citeauthor{DBLP:journals/pacmse/KhojahM0N24} investigated the use of ChatGPT (GPT-3.5) by professional software engineers in a week-long observational study~\cite{DBLP:journals/pacmse/KhojahM0N24}.
They found that most developers do not use the code generated by ChatGPT directly but instead use the output as a guide to implement their own solutions.
%Furthermore, they recommend that future research investigate the use of ChatGPT in non code-related SE tasks and for training and learning SE concepts.
\citeauthor{DBLP:conf/csee/AzanzaPIG24}'s case study found that LLMs could \enq{enhance personalized, instant onboarding support; however, relying on proprietary external LLMs poses significant data privacy risks}~\cite{DBLP:conf/csee/AzanzaPIG24}. 
 \citeauthor{DBLP:conf/icsa/JahicS24} found that most participants from 15 software companies they surveyed had already adopted AI (especially ChatGPT) for SE tasks; but cited copyright and privacy issues, as well as inconsistent or low-quality outputs, as barriers to adoption~\cite{DBLP:conf/icsa/JahicS24}. 
Retrospective studies that analyze data generated while developers use LLMs can provide additional insights into human-LLM interactions.
For example, researchers can employ data mining methods to build large-scale conversation datasets, such as the DevGPT dataset introduced by \citeauthor{DBLP:conf/msr/XiaoTHM24}~\cite{DBLP:conf/msr/XiaoTHM24}.
Conversations can then be analyzed using quantitative~\cite{DBLP:conf/msr/RabbiCZI24} and qualitative~\cite{DBLP:conf/msr/MohamedPP24} analysis methods.
 
\studytypesubsection{LLMs for New Software Engineering Tools (S6)}
\label{sec:llms-for-new-software-engineering-tools}

\studytypeparagraph{Description}

LLMs are being integrated into new tools that support software engineers in their daily tasks (e.g., to assist in code comprehension~\cite{DBLP:conf/chi/YanHWH24} and test case generation~\cite{DBLP:journals/tse/SchaferNET24}).
One way of integrating LLM-based tools into software engineers' workflows is using GenAI agents.
Unlike traditional LLM-based tools, these agents are capable of acting autonomously and proactively, are often tailored to meet specific user needs, and can interact with external environments~\cite{takerngsaksiri2024human,wiesinger2025agents}.
From an architectural perspective, GenAI agents can be implemented in various ways~\cite{wiesinger2025agents}.
However, they generally share three key components: (1) a reasoning mechanism that guides the LLM (often enabled by advanced prompt engineering), (2) a set of tools to interact with external systems (e.g., APIs or databases), and (3) a user communication interface that extends beyond traditional chat-based interactions~\cite{DBLP:conf/icsm/RichardsW24, DBLP:journals/tmlr/SumersYN024, DBLP:journals/corr/abs-2309-07870}.
Researchers can also test and compare different tool architectures to increase artifact quality and developer satisfaction.

\studytypeparagraph{Example(s)}

\citeauthor{DBLP:conf/chi/YanHWH24} proposed \emph{IVIE}, a tool integrated into the VS Code graphical interface that generates and explains code using LLMs~\cite{DBLP:conf/chi/YanHWH24}.
The authors focused more on the presentation, providing a user-friendly interface to interact with the LLM. 
\citeauthor{DBLP:journals/tse/SchaferNET24}~\cite{DBLP:journals/tse/SchaferNET24} presented a large-scale empirical evaluation on the effectiveness of LLMs for automated unit test generation.
They presented \emph{TestPilot}, a tool that implements an approach in which the LLM is provided with prompts that include the signature and implementation of a function under test, along with usage examples extracted from the documentation.
\citeauthor{DBLP:conf/icsm/RichardsW24} introduced a preliminary GenAI agent designed to assist developers in understanding source code by incorporating a reasoning component grounded in the theory of mind~\cite{DBLP:conf/icsm/RichardsW24}.
 
\studytypesubsection{Benchmarking LLMs for Software Engineering Tasks (S7)}
\label{sec:benchmarking-llms-for-software-engineering-tasks}

\studytypeparagraph{Description}

Benchmarking is the process of evaluating an LLM's performance on standardized tasks, using standardized metrics, on a standardized dataset. The LLM output is compared to a supposed ground truth from the benchmark dataset. Typical tasks include code generation, code summarization, code completion, and code repair, but also natural language processing tasks, such as anaphora resolution (i.e., the task of identifying the referring expression of a word or phrase occurring earlier in the text). %, interesting for subfields such a Requirements Engineering. 
Metrics may include general metrics for text generation, such as \emph{ROUGE}, \emph{BLEU}, or \emph{METEOR}~\cite{10.1145/3695988}, or task-specific metrics, such as \emph{CodeBLEU} for code generation.  Benchmarking requires high-quality reference datasets.

\studytypeparagraph{Example(s)}

In SE, benchmarking may include evaluating an LLM's ability to produce accurate and reliable outputs for a given input (usually a task description, possibly accompanied by data obtained from curated real-world projects or from synthetic SE-specific datasets). \emph{RepairBench}~\cite{silva2024repairbench}, for example, contains 574 buggy Java methods and their corresponding fixed versions, which can be used to evaluate the performance of LLMs in code repair tasks. It uses the \emph{Plausible@1} metric (i.e., the probability that the first generated patch passes all test cases) and the \emph{AST Match@1} metric (i.e., the probability that the abstract syntax tree of the first generated patch matches the ground truth patch).
\emph{SWE-Bench}~\cite{DBLP:conf/iclr/JimenezYWYPPN24} is a more generic benchmark that contains 2,294 SE Python tasks extracted from GitHub pull requests.
To score the LLM's task performance, the benchmark validates whether the generated patch successfully compiles and calculates the percentage of passed test cases. Meanwhile, \emph{HumanEval}~\cite{DBLP:journals/corr/abs-2107-03374} is often used to assess code generation. 
 
\subsection{Advantages and Challenges}
\label{sec:advantages-challenges}

Using LLMs in empirical SE research, whether as tools for researchers or for software engineers, offers several potential advantages but also raises fundamental challenges that cut across the study types described above.

\studytypeparagraph{Advantages}

The primary advantage of using LLMs for research tasks is \emph{speed, cost reduction, and scalability}. LLMs can annotate, judge, synthesize, and simulate faster and at lower cost than human researchers, with studies showing cost reductions of 50--96\% on various natural language tasks~\cite{DBLP:conf/emnlp/WangLXZZ21} (e.g., \citeauthor{DBLP:conf/chi/HeHDRH24} found that GPT-4 annotation required only two days and 122.08 USD compared to several weeks and 4,508 USD for a comparable MTurk pipeline~\cite{DBLP:conf/chi/HeHDRH24}). Similarly, augmenting or replacing human participants with LLM-generated virtual participants would reduce recruitment effort~\cite{DBLP:conf/vl/Madampe0HO24}. This efficiency unlocks scalability: qualitative research traditionally does not scale well to large samples, but LLMs allow researchers to analyze larger datasets~\cite{DBLP:journals/ase/BanoHZT24, barros2024largelanguagemodelqualitative, leça2024applicationsimplicationslargelanguage} and support larger judgment datasets.

LLMs can also \emph{automate} tasks such as coding qualitative data, assessing artifact quality, and generating summaries, reducing cognitive demands and resources required for qualitative research. Some studies suggest that LLMs may improve \emph{consistency}: ChatGPT's accuracy exceeded crowd workers by approximately 25\%, and LLMs can achieve higher inter-rater agreement than crowd workers and trained annotators~\cite{DBLP:journals/corr/abs-2303-15056}. However, no compelling evidence supports claims that LLMs are less biased than human annotators~\cite{DBLP:journals/corr/abs-2303-15056}.

LLMs could potentially provide \emph{access to otherwise-inaccessible research contexts}. \textit{If} virtual participants' behavior is sufficiently similar to human participants, LLMs could access underrepresented and hard-to-reach populations, strengthen generalizability and inclusiveness, impute missing data (see \synthesis), and enable research that is ethically problematic with real humans (e.g., questions that would force a human to relive past trauma do not harm an LLM).

Studying real-world LLM-based tools allows researchers to \emph{understand the state of practice}, uncovering usage patterns, adoption rates, and contextual factors, and \emph{generate hypotheses} about how LLMs affect developer productivity, collaboration, and decision-making. From an engineering perspective, developing LLM-based tools \emph{requires less task-specific engineering} than traditional SE approaches such as static analysis or symbolic execution, because a single model can handle diverse inputs without language-specific parsers or hand-crafted rules. Good benchmarks provide \emph{standardized evaluation and model comparison}, reducing research effort and fostering open science by providing common ground for sharing data and results. Benchmarks built for specific SE tasks can help identify LLM weaknesses and support optimization and fine-tuning.

\studytypeparagraph{Challenges}

The stochastic nature of LLM responses, where identical prompts may yield different outputs, \emph{undermines replicability} across all study types, complicating interpretation of experimental results and test-retest reliability. More broadly, \emph{reliability} is a persistent concern: conceptualizing an LLM-as-judge or LLM-as-annotator as a measurement instrument, researchers should expect validity, test-retest reliability, inter-rater reliability, minimal error, and measurement invariance~\cite{ralph2024teaching}. However, LLMs show significant variability depending on the dataset and task~\cite{DBLP:journals/corr/abs-2306-00176}, are sensitive to prompt variations~\cite{DBLP:journals/corr/abs-2304-11085} and option order~\cite{DBLP:conf/naacl/PezeshkpourH24}, can behave differently when reviewing their own outputs~\cite{NEURIPS2024_7f1f0218}, and can be unreliable for high-stakes labeling~\cite{DBLP:conf/chi/Wang0RMM24}. Crucially, \emph{reliability does not imply validity}: a reliable LLM might be reliably inaccurate~\cite{DBLP:conf/coling/ZhouC024}, and for tasks with no single correct answer, the statistical framework pushes outputs toward the most likely answer, which may not be the best one~\cite{DBLP:journals/corr/abs-2406-18403}. See \modelversion and \limitationsmitigations for reporting guidance on replicability and reliability.

The \emph{rapid evolution} of LLMs and LLM-based tools, combined with professionals rapidly learning how to employ them, complicates longitudinal comparisons and may quickly render study findings obsolete (e.g., GPT-4's accuracy on a simple mathematical task dropped from 84\% to 51\% within three months~\cite{DBLP:journals/corr/abs-2307-09009}).

The prevalence of \emph{proprietary tools and opaque training data} limits researchers' ability to assess and mitigate biases, and enables benchmark contamination~\cite{DBLP:journals/corr/abs-2410-16186}. LLMs exhibit \emph{bias} in multiple forms: tendencies to overestimate certain labels~\cite{DBLP:journals/corr/abs-2304-10145}, well-documented fairness issues~\cite{Gallegos2024BiasAF}, and the potential to reinforce prejudices. When simulating human participants, LLMs are \enq{likely to both misportray and flatten the representations of demographic groups}~\cite{DBLP:journals/corr/abs-2402-01908}. See \openllm and \limitationsmitigations for mitigation strategies.

\emph{Evaluation difficulty} is inherent in studying LLM-based tools and benchmarks. Benchmarks may lack construct validity~\cite{ralph2024teaching}, usually do not capture the full complexity of software engineering work~\cite{Chandra2025benchmarks}, and may lead to \emph{overconfidence} and overfitting~\cite{DBLP:conf/kdd/WanSJKCNSSWYABJ24}. LLMs are also susceptible to adversarial manipulation where semantics-preserving changes may deceive them into accepting flawed artifacts~\cite{DBLP:journals/corr/abs-2510-24367}. See \benchmarksmetrics for detailed guidance on benchmark design, including~\citeauthor{cao2025should}'s~\cite{cao2025should} guidelines for coding task benchmarks.

A fundamental challenge is \emph{philosophical and methodological incongruence} with qualitative research. In a recent open letter, 419 experienced qualitative researchers argued \enq{that analytic approaches such as reflexive thematic analysis are human research practices requiring a subjective, positioned, and reflexive researcher and therefore the use of GenAI in such approaches is not methodologically congruent}~\cite{jowsey2025reject}. LLMs lack the capacity for genuinely reflexive qualitative analysis because they operate on statistical prediction without understanding the meaning of the data being analyzed~\cite{jowsey2025reject, bano2023exploringqualitativeresearchusing, DBLP:journals/corr/abs-2510-18456}. Effective use requires structured prompts and careful human oversight~\cite{barros2024largelanguagemodelqualitative}; see \humanvalidation and \limitationsmitigations for guidance.

There is \emph{insufficient evidence} for the effectiveness of LLMs in most research roles. No compelling evidence exists that LLMs can accurately simulate human participants~\cite{schroeder2025llmspsychology} or reliably judge most relevant properties of SE artifacts; gathering such evidence for each specific usage may be quite difficult~\cite{DBLP:journals/ais/HardingDLL24}. LLMs may be better suited for augmenting rather than replacing human researchers~\cite{DBLP:conf/emnlp/WangLXZZ21, DBLP:conf/chi/HeHDRH24}, but even then, limited evidence exists that augmentation increases \emph{effectiveness} as well as \emph{efficiency}~\cite{schroeder2025llmspsychology}. When LLM outputs are incorrect, they can negatively \emph{affect human judgment}~\cite{DBLP:conf/www/HuangKA23a, panHumanCenteredDesignRecommendations2024}; see \humanvalidation for mitigation strategies.

Pooling multiple outputs for majority voting improves reliability but increases \emph{cost and environmental impact}~\cite{DBLP:journals/corr/abs-2304-11085, DBLP:conf/emnlp/WangLXZZ21}. While open models are available, the most capable ones require substantial hardware; relying on \emph{cloud-based APIs} introduces concerns related to data privacy, security, and replicability. See \limitationsmitigations for sustainability considerations.

Field study findings face \emph{generalizability} challenges because outcomes may be highly sensitive to individual differences, usage patterns, goals, and contexts. Field studies must be ``dependable''~\cite{Sullivan2011-ub} beyond traditional validity criteria, which complicates methodology; see \limitationsmitigations for a detailed threat taxonomy.
 
\section{Guidelines}
\label{sec:guidelines}

A primary goal of our guidelines is to \emph{enable reproducibility and replicability} of empirical SE studies involving LLMs. As repeating LLM-focused research to verify results lies somewhere between the ACM's definitions of reproducibility (different team, same research artifacts) and repeatability (different team, different research artifacts)~\cite{acm2020artifactreview} due to potential model changes and imperfect research artifacts, we follow \citeauthor{DBLP:journals/corr/abs-2510-25506}~\cite{DBLP:journals/corr/abs-2510-25506} and use the terms interchangeably. While previous guidelines regarding open science and empirical studies still apply, LLM-specific characteristics (e.g. inherent non-determinism~\cite{DBLP:conf/naacl/SongWLL25, YuanNondeterminismLLMInference}, opaque and proprietary models) present additional replicability challenges, which, in turn, demand new guidance.

Each guideline below begins with a brief \tldr summary, followed by its \textbf{rationale}, \textbf{recommendations}, \textbf{examples}, \textbf{benefits}, \textbf{challenges}, links to the relevant \textbf{study types}.
The \emph{rationale} articulates the underlying principle, i.e., why the guideline matters, while the \emph{recommendations} provide concrete, actionable practices.
\ifpaper
Table~\ref{tab:rationale-recommendations} provides a compact overview of this mapping.
\fi

The guidelines further contain \textbf{advice for reviewers}.
Broadly, reviewers should use our guidelines to help them interrogate the extent to which a manuscript's authors have done what was practically possible to improve reproducibility. We must neither accept research absent reasonable efforts to improve reproducibility, nor reject research for failing to obtain an impossible goal. We borrow this principle of ``reasonable efforts'' from the SIGSOFT Empirical Standards~\cite{ralph2021empiricalstandardssoftwareengineering}, where it applies to methodological rigor more generally. Of course, papers should acknowledge their limitations, but to determine whether these limitations are reasonable, reviewers should ask ``have the authors done what they could to minimize limitations?'' Our guidelines attempt to capture what ``reasonable efforts'' practically means for LLM-based studies.

To distinguish essential criteria from recommendations, our guidelines use two tiers.
A~\must criterion is a \emph{requirement}. Studies that intend to follow these
guidelines are expected to meet all \must criteria.
A~\should criterion represents a desired practice that strengthens a study's rigor or transparency.
However, there may be valid reasons to deviate in particular circumstances (e.g., resource constraints, inapplicability to a specific study context or type).
Nonetheless, studies that deviate from a \should criterion should briefly justify the
deviation and discuss its potential impact (e.g., on validity or reproducibility).
\ifpaper
Table~\ref{tab:guidelines} lists our eight guidelines.
Table~\ref{tab:guideline-matrix} provides an overview of how each guideline applies to each study type.
Cells marked \iconM indicate that the guideline is a \must requirement for the study type and cells marked \iconS indicate that the guideline \should be followed.
Cells marked -- indicate that the guideline is typically not applicable or not feasible for the study type.
\fi

The following sections indicate which information we expect researchers to report, and whether it should be in the \paper or \supplementarymaterial. Where a publication venue's page limits hinder reporting all expected elements in the \paper, it is better to report essential information in the \supplementarymaterial than not at all.
The \supplementarymaterial should be published according to the \emph{ACM SIGSOFT Open Science Policies}~\citep{daniel_graziotin_2024_10796477}.
 
\begin{table}[tb]
\caption{Overview of guidelines.}
\label{tab:guidelines}
\begin{tabular}{@{}lll@{}}
\toprule
\textbf{ID} & \textbf{Section} & \textbf{Guideline} \\
\midrule
G1 & \ref*{sec:declare-llm-usage-and-role} & \hyperref[sec:declare-llm-usage-and-role]{Declare LLM Usage and Role} \\
G2 & \ref*{sec:report-model-version-configuration-and-customizations} & \hyperref[sec:report-model-version-configuration-and-customizations]{Report Model Version, Configuration, and Customizations} \\
G3 & \ref*{sec:report-tool-architecture-beyond-models} & \hyperref[sec:report-tool-architecture-beyond-models]{Report Tool Architecture Beyond Models} \\
G4 & \ref*{sec:report-prompts-their-development-and-interaction-logs} & \hyperref[sec:report-prompts-their-development-and-interaction-logs]{Report Prompts, their Development, and Interaction Logs} \\
G5 & \ref*{sec:use-human-validation-for-llm-outputs} & \hyperref[sec:use-human-validation-for-llm-outputs]{Use Human Validation for LLM Outputs} \\
G6 & \ref*{sec:use-an-open-llm-as-a-baseline} & \hyperref[sec:use-an-open-llm-as-a-baseline]{Use an Open LLM as a Baseline} \\
G7 & \ref*{sec:use-suitable-baselines-benchmarks-and-metrics} & \hyperref[sec:use-suitable-baselines-benchmarks-and-metrics]{Use Suitable Baselines, Benchmarks, and Metrics} \\
G8 & \ref*{sec:report-limitations-and-mitigations} & \hyperref[sec:report-limitations-and-mitigations]{Report Limitations and Mitigations} \\
\bottomrule
\end{tabular}
\end{table}

\providecommand{\rot}[1]{\rotatebox{60}{\textbf{#1}}}

\begin{table}[tb]
\caption{Applicability of guidelines to study types.
\iconM~=~the guideline's core recommendations \must be followed for this study type,
\iconS~=~\should be followed,
--~=~are not directly applicable.
Study type IDs: S1~=~\annotators, S2~=~\judges, S3~=~\synthesis, S4~=~\subjects, S5~=~\llmusage, S6~=~\newtools, S7~=~\benchmarkingtasks.
Each guideline's study-type-specific guidance is detailed in the corresponding subsection.}
\label{tab:guideline-matrix}
\footnotesize
\setlength{\tabcolsep}{4pt}
\renewcommand{\arraystretch}{1.15}
\begin{tabular}{@{} l c c c c c c c @{}}
\toprule
& \rot{\hyperref[sec:llms-as-annotators]{S1}} & \rot{\hyperref[sec:llms-as-judges]{S2}} & \rot{\hyperref[sec:llms-for-synthesis]{S3}} & \rot{\hyperref[sec:llms-as-subjects]{S4}} & \rot{\hyperref[sec:studying-llm-usage-in-software-engineering]{S5}} & \rot{\hyperref[sec:llms-for-new-software-engineering-tools]{S6}} & \rot{\hyperref[sec:benchmarking-llms-for-software-engineering-tasks]{S7}} \\
\midrule
\hyperref[sec:declare-llm-usage-and-role]{G1: Declare Usage}\hspace{12pt}
  & \iconM & \iconM & \iconM
  & \iconM & \iconM & \iconM
  & \iconM \\
\hyperref[sec:report-model-version-configuration-and-customizations]{G2: Model Version}
  & \iconM & \iconM & \iconM
  & \iconM & \iconM & \iconM
  & \iconM \\
\hyperref[sec:report-tool-architecture-beyond-models]{G3: Architecture}
  & \iconS & \iconS & \iconS
  & \iconS & \iconS & \iconM
  & \iconS \\
\hyperref[sec:report-prompts-their-development-and-interaction-logs]{G4: Prompts}
  & \iconM & \iconM & \iconM
  & \iconM & \iconM & \iconM
  & \iconM \\
\hyperref[sec:use-human-validation-for-llm-outputs]{G5: Human Valid.}
  & \iconS & \iconS & \iconS
  & \iconS & \iconS & \iconS
  & \iconS \\
\hyperref[sec:use-an-open-llm-as-a-baseline]{G6: Open LLM}
  & \iconS & \iconS & \iconS
  & -- & -- & \iconS
  & \iconM \\
\hyperref[sec:use-suitable-baselines-benchmarks-and-metrics]{G7: Benchmarks}
  & \iconM & \iconS & \iconS
  & \iconS & \iconS & \iconS
  & \iconM \\
\hyperref[sec:report-limitations-and-mitigations]{G8: Limitations}
  & \iconM & \iconM & \iconM
  & \iconM & \iconM & \iconM
  & \iconM \\
\bottomrule
\end{tabular}
\end{table} 
\subsection{Declare LLM Usage and Role (G1)}
\label{sec:declare-llm-usage-and-role}

\vspace{0.5\baselineskip}
\begin{framed}
\tldr: Researchers \must disclose any use of LLMs to support empirical studies in their \paper, specifying which LLM was used, how it was used, and where in the research process it was employed. They \should report the exact purpose, the tasks that were automated, and the expected benefits in the \paper. When the LLM plays a central role, the declaration \should be prominent and detailed in the methodology section; for tangential uses, a brief statement in the methodology or acknowledgments suffices. \end{framed}

\guidelinesubsubsection{Rationale}

Transparency about LLM involvement is a prerequisite for informed assessment of a study's scope, limitations, and potential biases.
Without explicit disclosure, readers cannot evaluate how the LLM's characteristics may have influenced the research process or its outcomes.

\guidelinesubsubsection{Recommendations}

When conducting any kind of empirical study involving LLMs, researchers \must clearly declare that an LLM was used (see \scope for what we consider relevant research support).
This \should be done in a suitable section of the \paper, for example, in the introduction or research methods section.
For authoring scientific articles, this transparency is, for example, required by the ACM Policy on Authorship: \enq{The use of generative AI tools and technologies to create content is permitted but must be fully disclosed in the Work}~\cite{ACM2023}.

Beyond generic declarations, researchers \should report the exact purpose of using an LLM in a study, the tasks it was used to automate, and the expected benefits in the \paper.
A sufficient declaration specifies not only \emph{that} an LLM was used, but also \emph{which} LLM (name and version), \emph{how} it was used (e.g., as an annotator, code generator, or judge), and \emph{where} in the research process it was employed (e.g., data collection, analysis, or synthesis).

When the LLM plays a central role in the study (e.g., as the main tool being evaluated or as a core component of the research method), the declaration \should be prominent and detailed, appearing in the methodology section with cross-references to the specific guidelines that apply (e.g., Sections \modelversion, \toolarchitecture, and \prompts).
When the LLM's role is more tangential (e.g., used for a single preprocessing step), a brief but explicit statement in the methodology or acknowledgments section is sufficient.
In either case, the disclosure \must be specific enough for readers to assess how the LLM's involvement may affect the study's validity and reproducibility.
% Related paper:
% Artificial intelligence‑assisted academic writing: recommendations for ethical use
% https://doi.org/10.1186/s41077-025-00350-6

\guidelinesubsubsection{Example(s)}

The \emph{ACM Policy on Authorship}~\cite{ACM2023} suggests disclosing GenAI usage in the acknowledgments section of the \paper, advising to ``err on the side of caution, and include a disclosure in the acknowledgments section of the Work'' when uncertain about the need~\cite{ACM2023}.
For double-blind review, researchers can add a temporary ``AI Disclosure'' section where the acknowledgments would appear.
An example of an LLM disclosure beyond writing support can be found in a recent paper by \citeauthor{DBLP:conf/re/LubosFTGMEL24}~\cite{DBLP:conf/re/LubosFTGMEL24}, in which they write in the methodology section:

\begin{quote}
\it
``We conducted an LLM-based evaluation of requirements utilizing the Llama 2 language model with 70 billion parameters, fine-tuned to complete chat responses...''
\end{quote}

\guidelinesubsubsection{Benefits}

Transparency in the use of LLMs helps other researchers understand the context and scope of the study, facilitating better interpretation and comparison of the results.
Beyond this declaration, we recommend researchers to be explicit about the LLM version they used (see \modelversion) and the LLM's exact role.

\guidelinesubsubsection{Challenges}

Declaring LLM usage requires only a brief statement and no additional experiments, making compliance straightforward.
One challenge might be authors' reluctance to disclose LLM usage for valid use cases, because they fear that AI-generated content makes reviewers think that the authors' work is less original.
In fact, there is evidence suggesting that AI disclosure can negatively affect trust in authors~\cite{SCHILKE2025104405}.
However, the \emph{ACM Policy on Authorship} is very clear in that any use of GenAI tools to create content \must be disclosed.
Our guidelines focus on research support beyond proof-reading and writing support (see \scope), but the threshold of what must be declared continues to evolve as organizations such as the ACM update their authorship policies.

\guidelinesubsubsection{Study Types}

Researchers \must follow this guideline for all study types.
The specific focus of the declaration varies by study type.
For \annotators, \judges, \synthesis, and \subjects, researchers \must declare the specific role assigned to the LLM (e.g., annotator, judge, synthesizer, or simulated participant).
For \llmusage, researchers \must clarify which LLM(s) the observed participants used and under which conditions.
For \newtools, researchers \must declare the LLM's role within the tool architecture and its contribution to the tool's functionality.
For \benchmarkingtasks, researchers \must declare which LLMs were benchmarked and for which tasks.

\guidelinesubsubsection{Advice for Reviewers}

The most common problem with disclosure is incompleteness or vagueness about how the LLM was used. If the paper says ``we used LLM X to help with task Y'' without specifying how, reviewers should request clarification. Such requests are typically minor revisions unless the missing details may reveal methodological problems.

Using an LLM as a copyeditor becomes problematic when authors do not or cannot take responsibility for the resulting text. If a reviewer finds clear evidence that LLM-generated text was not carefully reviewed (e.g., text like ``As a Large Language Model, I...''), this may warrant rejection. If a reviewer suspects an author \textit{cannot} take responsibility for AI-generated text, they should raise the concern with their editor or program chair. Due process requires that authors not be accused of misconduct without clear evidence.

If undisclosed LLM use is suspected, the reviewer should similarly consult their editor or program chair. When the evidence is conclusive, the key question is the degree to which undisclosed use affects the study's contribution, ranging from negligible (e.g., word choice in a single sentence) to severe (e.g., generating ostensibly empirical data or statistical analyses).

\guidelinesubsubsection{See Also}
\begin{itemize}[label=$\rightarrow$]
    \item Section~\modelversion: model version and configuration details.
    \item Section~\toolarchitecture: system-level role documentation.
    \item Section~\prompts: prompt documentation and interaction logs.
\end{itemize} 
\subsection{Report Model Version, Configuration, and Customizations (G2)}
\label{sec:report-model-version-configuration-and-customizations}

\vspace{0.5\baselineskip}
\begin{framed}
\tldr: Researchers \must report the exact LLM model or tool version, configuration, and experiment date in the \paper. When using quantized models, researchers \should report the quantization level and method. For fine-tuned models, they \must describe the fine-tuning goal, dataset, and procedure in the \paper. Researchers \should include default parameters, explain model choices, compare base- and fine-tuned model using suitable metrics and benchmarks, and share fine-tuning data and weights as \supplementarymaterial (or alternatively justify in the \paper why they cannot share them). \end{framed}

\guidelinesubsubsection{Rationale}

LLMs and LLM-based tools are frequently updated, and configuration parameters such as temperature or seed values affect content generation. 
This guideline focuses on documenting the \emph{model-specific} aspects of empirical studies involving LLMs, concentrating on the models themselves, their version, configuration parameters, and customizations (e.g., fine-tuning). While the \toolarchitecture section addresses system-level integration, the information outlined here is always essential for reproducibility whenever an LLM is involved.

\guidelinesubsubsection{Recommendations}

Researchers \must document in the \paper which model or tool version they used in their study, along with the date when the experiments were carried out and the configured parameters that affect output generation.
Since default values might change over time, researchers \should always report all configuration values, even if they used the defaults.
Checksums and fingerprints \should be reported since they identify specific versions and configurations.
Depending on the study context, other properties such as the context window size (number of tokens) \should be reported.
When using quantized models, researchers \should report the quantization level (e.g., 4-bit, 8-bit) and method (e.g., GPTQ or AWQ), as different quantization approaches produce different outputs, affecting both output quality and reproducibility.
Researchers \should motivate in the \paper why they selected certain models, versions, and configurations.
Reasons may be monetary, technical, or methodological (e.g., planned comparison to previous work).
Depending on the specific study context, additional information regarding the experiment or tool architecture \should be reported.

A common customization approach for existing LLMs is fine-tuning.
If a model was fine-tuned, researchers \must describe the fine-tuning goal (e.g., improving the performance for a specific task), the fine-tuning procedure (e.g., full fine-tuning vs. Low-Rank Adaptation (LoRA), selected hyperparameters, loss function, learning rate, batch size, etc.), and the fine-tuning dataset (e.g., data sources, the preprocessing pipeline, dataset size) in the \paper.
Researchers \should either share the fine-tuning dataset as part of the \supplementarymaterial or explain in the \paper why the data cannot be shared (e.g., because it contains confidential or personal data that could not be anonymized).
The same applies to the fine-tuned model weights.
Suitable benchmarks and metrics \should be used to compare the base model with the fine-tuned model.

In summary, our recommendation is to report:
\begin{enumerate*}[label=(\arabic*)]
\item Model/tool name and version (\must in \paper);
\item All relevant configured parameters that affect output generation (\must in \paper);
\item Default values of all available parameters (\should);
\item Checksum/fingerprint of used model version and configuration (\should);
\item Additional properties such as context window size (\should);
\item Model quantization level and method, if applicable (\should).
\end{enumerate*}

For fine-tuned models, additional recommendations apply:
\begin{enumerate*}[label=(\arabic*)]
\item Fine-tuning goal (\must in \paper);
\item Fine-tuning dataset creation and characterization (\must in \paper);
\item Fine-tuning parameters and procedure (\must in \paper);
\item Fine-tuning dataset and fine-tuned model weights (\should);
\item Validation metrics and benchmarks (\should).
\end{enumerate*}

Commercial models (e.g., GPT-5) or LLM-based tools (e.g., ChatGPT) might not give researchers access to all required information.
Our suggestion is to report what is available and openly acknowledge limitations that hinder reproducibility.

\guidelinesubsubsection{Example(s)}

Based on the documentation that OpenAI and Azure provide~\cite{OpenAI25, Azure25}, researchers might, for example, report:

\begin{quote}
\it
 ``We integrated a  \texttt{gpt-4} model in version \texttt{0125-Preview} via the Azure OpenAI Service, and configured it with a temperature of 0.7, top\_p set to 0.8, a maximum token length of 512, and the  seed value \texttt{23487}.
 We ran our experiment on 10th January 2025. The system fingerprint was \texttt{fp\_6b68a8204b}.
\end{quote}

\citeauthor{DBLP:conf/icse/KangYY23} provide a similar statement in their paper on exploring LLM-based bug reproduction~\cite{DBLP:conf/icse/KangYY23}:

\begin{quote}
\it
``We access OpenAI Codex via its closed beta API, using the code-davinci-002 model. For Codex, we set the temperature to 0.7, and the maximum number of tokens to 256.''
\end{quote}

Our guidelines additionally suggest to report a checksum/fingerprint and exact dates, but otherwise this example is close to our recommendations. 

\citeauthor{DBLP:conf/icsa/DharVV24}~\cite{DBLP:conf/icsa/DharVV24} assessed whether LLMs can generate architectural design decisions, detailing the system architecture and the LLM's role within it.
They provide information on the fine-tuning approach and datasets, including the source of architectural decision records, preprocessing methods, and data selection criteria.

For self-hosted models, the \supplementarymaterial can become a true replication package. For example, for models provisioned using \href{https://ollama.com/library/}{ollama}, one can report the specific tag and checksum, e.g., \emph{``llama3.3, tag 70b-instruct-q8\_0, checksum d5b5e1b84868.''}
Given suitable hardware, running the model is then as easy as executing the following command:\\
\texttt{ollama run llama3.3:70b-instruct-q8\_0}

\guidelinesubsubsection{Benefits}

Reporting this information is a prerequisite for the verification, reproduction, and replication of LLM-based studies. While LLMs are inherently non-deterministic, this cannot excuse dismissing reproducibility. Although exact reproducibility is hard to achieve, reporting the information outlined in this guideline helps researchers come as close as possible to that standard.
%However, that information alone is generally not sufficient.

\guidelinesubsubsection{Challenges}

Different model providers and modes of operating the models allow for varying degrees of information.
For example, OpenAI provides a model version and a system fingerprint describing the backend configuration, which can also influence the output.
However, in fact, the fingerprint is intended only to detect changes in the model or its configuration; one cannot go back to a certain fingerprint.
As a beta feature, OpenAI lets users set a seed parameter to receive ``(mostly) consistent output''~\cite{OpenAI23}.
However, the seed value does not allow for full reproducibility and the fingerprint changes frequently. 
Although, as motivated above, open models significantly simplify re-running experiments, they also come with challenges in terms of reproducibility, as generated outputs can be inconsistent despite setting the temperature to 0 and using a seed value (see \href{https://github.com/ollama/ollama/issues/5321}{GitHub issue for Llama3}).
Setting the temperature to 0 configures greedy decoding (always selecting the most probable next token), which minimizes output variability but can degrade quality by producing repetitive text and missing higher-quality responses~\cite{holtzman2020curious}.

Even with a temperature of 0, full determinism is rarely guaranteed: floating-point arithmetic on GPUs causes slight numerical differences that cascade into divergent token selections~\cite{YuanNondeterminismLLMInference}, Sparse Mixture-of-Experts routing amplifies this effect~\cite{Chann2023}, silent backend changes in commercial APIs produce different outputs over time~\cite{DBLP:journals/corr/abs-2307-09009}, and even self-hosted open models with identical settings do not always yield consistent outputs~\cite{DBLP:journals/corr/abs-2510-25506, DBLP:conf/naacl/SongWLL25}.
Researchers should therefore not treat a temperature of 0 as a guarantee of reproducibility, but as one measure among several, including fixed seed values~\cite{OpenAI23}, system fingerprints, and archiving of raw outputs.
When a temperature of 0 is chosen primarily for reproducibility, this motivation \should be stated explicitly, along with an acknowledgment of its potential impact on output quality.
%\comment{Should we also address how API rate limits and usage constraints can affect experiments?}

\guidelinesubsubsection{Study Types}

This guideline \must be followed for all study types for which the researcher has access to (parts of) the model's configuration.
They \must always report the configuration that is visible to them, acknowledging the reproducibility challenges of commercial tools and models that are offered as-a-service.
Depending on the specific study type, researchers \should provide additional information on the architecture of a tool they built (see \toolarchitecture), prompts and interactions logs (see \prompts), and specific limitations and mitigations (see \limitationsmitigations).

For example, when \llmusage by focusing on commercial tools such as ChatGPT or GitHub Copilot, researchers \must be as specific as possible in describing their study setup.
The configured model name, version, and the date when the experiment was conducted \must always be reported.
In those cases, reporting other aspects, such as prompts and interaction logs, is essential.

For \annotators, \judges, and \synthesis, researchers \must report the model configuration used for the respective annotation, judging, or synthesis tasks, including temperature and other sampling parameters that affect output variability.
For \subjects, researchers \must report any persona-related configuration settings and parameters that shape the simulated behavior.
For \newtools, researchers \must report the configuration for each model integrated in the tool's architecture, including any model-specific parameter choices.
For \benchmarkingtasks, researchers \must report the configuration for all benchmarked models to enable fair cross-model comparisons.

\guidelinesubsubsection{Advice for Reviewers}

Missing version, configuration, or parameter information is typically a minor revision request. Before concluding that information is absent, reviewers should check appendices and supplementary materials, as details are sometimes reported there rather than in the main text. Rejection over missing details is rarely warranted unless the omissions obscure deeper methodological problems.

\guidelinesubsubsection{See Also}
\begin{itemize}[label=$\rightarrow$]
    \item Section~\toolarchitecture: system-level architecture beyond the model.
    \item Section~\prompts: prompt reporting and development.
    \item Section~\benchmarksmetrics: benchmarks for comparing base and fine-tuned models.
    \item Section~\openllm: open models for self-hosted reproducibility.
    \item Section~\limitationsmitigations: reporting reproducibility limitations.
\end{itemize} 
\subsection{Report Tool Architecture Beyond Models (G3)}
\label{sec:report-tool-architecture-beyond-models}

\vspace{0.5\baselineskip}
\begin{framed}
\tldr: Researchers \must describe the full architecture of LLM-based tools that they develop in their \paper. This includes the role of the LLM, interactions with other components, and the overall system behavior. If autonomous agents are used, researchers \must specify agent roles, reasoning frameworks, and communication flows. Hosting, hardware setup, and latency implications \must be reported. For tools using retrieval or augmentation methods, data sources, integration mechanisms, and update and versioning strategies \must be described. For ensemble architectures, the coordination logic between models \must be explained in the \paper. Researchers \should include architectural diagrams and justify design decisions. Non-disclosed confidential or proprietary components \must be acknowledged as reproducibility limitations. \end{framed}

\guidelinesubsubsection{Rationale}

LLM-based tools often have \emph{complex software layers} around the model(s) that pre-process data, prepare prompts, filter user requests, or post-process responses.
For example, ChatGPT and GitHub Copilot use the same underlying models, but their outputs differ significantly because Copilot automatically adds project context; researchers can also build tools using models directly via APIs.
The infrastructure and business logic around the bare model can significantly contribute to the performance of a tool for a given task, for instance, RAG pipelines, output parsers, and retry logic all shape the final output independently of the model itself.

This section addresses the \emph{system-level} aspects of LLM-based tools that researchers develop, complementing \modelversion, which focuses on model-specific details.
While standalone LLMs require limited architectural documentation, a more detailed description is required for study setups and tool architectures that integrate LLMs with other components to create more complex systems.
This section provides guidelines for documenting these broader architectures.

\guidelinesubsubsection{Recommendations}

%\paragraph{General Requirements:}

Researchers \must clearly describe the tool architecture and what exactly the LLM (or ensemble of LLMs) contributes to the tool or method presented in a research paper.
Researchers \should provide a high-level architectural diagram to improve transparency.
To improve clarity, researchers \should explain design decisions, particularly regarding how the models were hosted and accessed (API-based, self-hosted, etc.) and which retrieval mechanisms were implemented (keyword search, semantic similarity matching, rule-based extraction, etc.).
Researchers \mustnot omit critical architectural details that could affect reproducibility, such as dependencies on proprietary tools that influence tool behavior.
Especially for \emph{time-sensitive measurements}, the description of the hosting environment is central, as it can significantly impact the results.
Researchers \must clarify whether local infrastructure or cloud services were used, including detailed infrastructure specifications and latency considerations.

\paragraph{Agent-Based and Complex Systems:}
If an LLM is used as a \emph{standalone system}, for example, by sending prompts directly to a GPT-4o model via the OpenAI API without pre-processing the prompts or post-processing the responses, a brief explanation of this approach is usually sufficient.
However, if LLMs are integrated into more \emph{complex systems} with pre-processing, retrieval mechanisms, or autonomous agents, researchers \must provide a detailed description of the system architecture in the \paper.
Aspects to consider are how the LLM interacts with other components such as databases, external APIs, and frameworks.
If the LLM is part of an \emph{agent-based system} that autonomously plans or executes tasks, researchers \must describe its exact architecture, including the agents' roles (e.g., planner, executor, coordinator), whether it is a single-agent or multi-agent system, how it interacts with external tools and users, and the reasoning framework used (e.g., chain-of-thought, self-reflection, multi-turn dialogue).
Any agent configuration via context files (e.g., \texttt{AGENTS.md}) \should also be documented as part of the system architecture (see \prompts for detailed reporting requirements).
For agent-based systems that use external tools (e.g., Claude Code and its \href{https://docs.anthropic.com/en/docs/claude-code/sub-agents}{subagents}), researchers \must discuss \emph{agent behavior traceability} by clearly delineating three distinct components: (1) \emph{LLM Input/Output}: Internal deliberation, planning, and interpretation performed by the model. (2) \emph{External tool calls}: Specific invocations of APIs, databases, file systems, or other external services that the agent explicitly triggers (e.g., via the \href{https://modelcontextprotocol.io/}{Model Context Protocol (MCP)}). (3) \emph{User or system interactions}: Human-in-the-loop feedback, environment responses, or multi-agent communication.
Researchers \should report complete execution traces that show the sequence and causality between these components, including inputs and outputs of all external tool calls.
This traceability is essential for determining whether task success is attributable to the LLM's output and tool-calling capabilities, the external tools' functionality, or their interaction patterns.
When full tool call logging is not feasible due to privacy or proprietary constraints, researchers \should provide representative examples or anonymized traces that demonstrate the agent's decision-making process.
%Researchers \mustnot present an agent-based system without detailing how it makes decisions and executes tasks.

\paragraph{Retrieval, Augmentation, and Ensembles:}
If \emph{retrieval or augmentation methods} were used (e.g., retrieval-augmented generation (RAG), rule-based retrieval, structured query generation, or hybrid approaches), researchers \must describe how external data is retrieved, stored, and integrated into the LLM's responses.
This includes specifying the type of storage or database used (e.g., vector databases, relational databases, knowledge graphs) and how the retrieved information is selected and used.
Stored data used for context augmentation \must be reported, including details on data preprocessing, versioning, and update frequency.
If this data is not confidential, an anonymized snapshot of the data used for context augmentation \should be made available as \supplementarymaterial.
For \emph{ensemble models}, in addition to following the \modelversion guideline for each model, the researchers \must describe the architecture that connects the models.
The \paper \must at least contain a high-level description, and details can be reported in the \supplementarymaterial.
Aspects to consider include documenting the logic that determines which model handles which input, the interaction between models, and the architecture for combining outputs (e.g., majority voting, weighted averaging, sequential processing).

\guidelinesubsubsection{Example(s)}

\citeauthor{DBLP:journals/tse/SchaferNET24}~\cite{DBLP:journals/tse/SchaferNET24} evaluated LLMs for automated unit test generation, providing a comprehensive description of the system architecture including code parsing, prompt formulation, LLM interaction, and test suite integration.
They also detail the datasets used, including sources, selection criteria, and preprocessing steps.

A second example is \citeauthor{DBLP:conf/chi/YanHWH24}'s \emph{IVIE} tool~\cite{DBLP:conf/chi/YanHWH24}, which integrates LLMs into the VS Code interface.
The authors document the tool architecture, detailing the IDE integration, context extraction from code editors, and the formatting pipeline for LLM-generated explanations.
This documentation illustrates how architectural components beyond the core LLM affect the overall tool performance and user experience.

\guidelinesubsubsection{Benefits}

The software layers around bare LLMs significantly impact tool performance and hence need detailed documentation.
Documenting the architecture and supplemental data of LLM-based systems enhances reproducibility and transparency~\cite{DBLP:journals/software/LuZXXW24}, enabling experiment replication, result validation, and cross-study comparison.

\guidelinesubsubsection{Challenges}

Researchers face challenges in documenting LLM-based architectures, including proprietary APIs and dependencies that restrict disclosure, managing large-scale retrieval databases, and ensuring efficient query execution.
They must also balance transparency with data privacy concerns, adapt to the evolving nature of LLM integrations, and, depending on the context, handle the complexity of multi-agent interactions and decision-making logic, all of which can impact reproducibility and system clarity.

\guidelinesubsubsection{Study Types}

This guideline \must be followed for all studies that involve tools with system-level components beyond bare LLMs, from lightweight wrappers that pre-process user input or post-process model outputs, to systems employing retrieval-augmented methods or complex agent-based architectures.

For \newtools, this guideline is of primary importance: researchers \must describe the tool's full architecture, including the role of each LLM, interactions with other components, and the overall system behavior.
For \benchmarkingtasks, researchers \should describe the evaluation harness and infrastructure if it goes beyond bare model API calls (e.g., custom sandboxing, orchestration layers, or post-processing pipelines).
For \llmusage, researchers \should describe the tool architecture of the studied tool to the extent that it is accessible, as architectural details may influence observed usage patterns.
For \annotators, \judges, \synthesis, and \subjects, if the research setup involves a custom pipeline (e.g., retrieval-augmented generation for annotation or chained prompts for synthesis), the architecture \should also be reported.

%% Changed: Fixed "not the only supplement" → "not only the supplement".
%% Suggestion: Consider restructuring the proprietary-components paragraph into clearer tiers: (1) evaluating opaque artifacts is a valid scientific contribution, (2) inscrutability weakens the contribution, (3) scientific instruments must be fully disclosed.
\guidelinesubsubsection{Advice for Reviewers}

As with other guidelines, missing architectural information is typically a minor revision request unless so much is missing that methodological rigor cannot be assessed. Reviewers may ask authors to move key details from supplements into the paper body to ensure the main text is self-contained.

Regarding proprietary or confidential components, three principles apply: (1) empirically evaluating an opaque artifact is a valid scientific contribution; (2) the less available and inspectable the artifact, the weaker the contribution; (3) \textit{scientific instruments} including tools, metrics, scales, and experimental materials,  must be fully disclosed, even when the object under study cannot be. 
\subsection{Report Prompts, their Development, and Interaction Logs (G4)}
\label{sec:report-prompts-their-development-and-interaction-logs}

\vspace{0.5\baselineskip}
\begin{framed}
\tldr: Researchers \must publish all prompts or, when using templates, prompt templates with representative instances, including their structure, content, formatting, and dynamic components, as \supplementarymaterial. If full prompt disclosure is not feasible, for example, due to privacy or confidentiality concerns, summaries or examples \should be provided. Prompt development strategies (e.g., zero-shot, few-shot), rationale, and selection process \must be described in the \paper. When prompts are long or complex, input handling and token optimization strategies \must be documented. For dynamically generated or user-authored prompts, generation and collection processes \must be reported. Prompt reuse across models and configurations \must be specified. Researchers \should report prompt revisions and pilot testing insights. For agentic systems, developed plans \should be reported as \supplementarymaterial and interaction logs \should be generalized to include human-agent interaction traces as \supplementarymaterial. All context files used to configure AI agent behavior \should be reported as \supplementarymaterial. To address model non-determinism and ensure reproducibility, especially when targeting SaaS-based commercial tools, full interaction logs (prompts and responses) \should be included as \supplementarymaterial if privacy and confidentiality can be ensured. \end{framed}

\guidelinesubsubsection{Rationale}

Prompts are critical for any study involving LLMs~\cite{DBLP:conf/iclr/Sclar0TS24}.
A \emph{prompt} is a concrete input to an LLM that guides its output~\cite{DBLP:journals/corr/abs-2406-06608}.
Depending on the task, a prompt may include instructions (e.g., \enq{classify the following bug report}), context (e.g., a taxonomy of bug categories, related source code), input data (e.g., the bug report text), and output format specifications (e.g., \enq{respond as JSON with fields `category' and `justification'}), with outputs ranging from unstructured text to structured formats such as JSON.\footnote{\url{https://www.promptingguide.ai/introduction/elements}}
A \emph{prompt template} is a parameterized structure containing static elements (e.g., instructions, output format specifications) and placeholders for dynamic content (e.g., source code under analysis) that are filled in at runtime to construct concrete prompts~\cite{DBLP:journals/corr/abs-2406-06608}.
In automated studies, researchers typically design prompt templates from which individual prompts are then instantiated.
Prompts significantly influence the quality of the output of a model, and understanding how exactly they were formatted and integrated into an LLM-based study is essential to ensure transparency, verifiability, and reproducibility.
\citeauthor{DBLP:conf/iclr/Sclar0TS24} question the methodological validity of comparing models with \enq{an arbitrarily chosen, fixed prompt format}, because their research has shown that the performance of different prompt formats only weakly correlates between different models~\cite{DBLP:conf/iclr/Sclar0TS24}.
We remind the reader that these guidelines do not apply when LLMs are used solely for language polishing, paraphrasing, translation, tone or style adaptation, or layout edits (see \scope).
This implies that our guidelines do not recommend that researcher disclose the prompts they have used for such tasks.

\guidelinesubsubsection{Recommendations}

\paragraph{Prompt Reporting:}
Researchers \must report all prompts that were used for an empirical study, including all instructions, context, input data, and output indicators.
The only exception to this is if transcripts might identify anonymous participants or reveal personal or confidential information, for example, when working with industry partners.
When prompt templates are used, researchers \must report them alongside representative prompt instances, specifying which parts are static and which are dynamically filled.
Moreover, researchers \should specify the exact formatting of prompts, including how code snippets were enclosed (e.g., markdown-style code blocks, triple backticks), whether whitespace was preserved, and how other artifacts such as error messages and stack traces were presented.

\paragraph{Prompt Development:}
Researchers \must also explain in the \paper how they developed the prompts and why they decided to follow certain prompting strategies.
If prompts from the early phases of a research project are unavailable, they \must at least summarize the prompt evolution.
However, given that prompts can be stored as plain text, there are established SE techniques such as version control systems that can be used to collect the required provenance information.

Prompt development is often iterative, involving collaboration between human researchers and AI tools.
Researchers \should report any instances in which LLMs were used to suggest prompt refinements and how these suggestions were incorporated.
Prompts may need revision in response to failure cases, and pilot testing is vital to ensure reliable results.
If such testing was conducted, researchers \should summarize key insights, including how different prompt variations affected output quality and which criteria were used to finalize the prompt design.
%Iterative changes based on human feedback and pilot testing results \should be included in the documentation.
A prompt changelog can track prompt evolution, including key revisions and reasons for changes (e.g., v1.0: initial prompt; v2.0: incorporated examples of ideal responses).
Since prompt effectiveness varies between models and model versions, researchers \must make clear which prompts were used for which models in which versions and with which configuration.

\paragraph{Prompting Strategy and Input Handling:}
Researchers \must specify whether zero-shot, one-shot, or few-shot prompting was used.
For few-shot prompts, the examples provided to the model \should be clearly outlined, along with the rationale for selecting them.
If multiple versions of a prompt were tested, researchers \should describe how these variations were evaluated and how the final design was chosen.

When dealing with extensive or complex prompt context, researchers \must describe the strategies they used to handle input length constraints.
Approaches might include truncating, summarizing, or splitting prompts into multiple parts.
Token optimization measures, such as simplifying code formatting or removing unnecessary comments, \must also be documented if applied.

\paragraph{Dynamic and User-Authored Prompts:}
In cases where prompts are generated dynamically, such as through preprocessing, template structures, or retrieval-augmented generation (RAG), the process \must be thoroughly documented.
This includes explaining any automated algorithms or rules that influenced prompt generation.
For studies involving human participants who create or modify prompts, researchers \must describe how these prompts were collected and analyzed.

\paragraph{Supplementary Material and Anonymization:}
To ensure complete reproducibility, researchers \must make all prompts and prompt variations publicly available as part of their \supplementarymaterial.
If the complete set of prompts cannot be included in the \paper, researchers \should provide summaries and representative examples.
This also applies if complete disclosure is not possible due to privacy or confidentiality concerns.
For prompts containing sensitive information, researchers \must: (1) anonymize personal identifiers, (ii) replace proprietary code with placeholders, and (iii) clearly highlight modified sections.

\paragraph{Interaction Logs:}
When trying to verify results, even with the exact same prompts, decoding strategies, and parameters, LLMs can still behave non-deterministically.
Non-determinism can arise from batching, input preprocessing, and floating point arithmetic on GPUs~\cite{Chann2023}.
Thus, in order to enable other researchers to verify the conclusions that researchers have drawn from LLM interactions, researchers \should report the full interaction logs (prompts and responses) as part of their \supplementarymaterial.
Reporting this is especially important for studies targeting commercial software-as-a-service (SaaS) solutions such as ChatGPT.
%or novel tools that integrate LLMs via cloud APIs where there is even less guarantee of reproducing the state of the LLM-powered system at a later point by a reader of the study who wants to verify, reproduce, or replicate it. 
The rationale for this is similar to the rationale for reporting interview transcripts in qualitative research.
Just as a human participant might give different answers to the same question asked two months apart, the responses from tools such as ChatGPT can also vary over time.
%Therefore, keeping a record of the actual conversation is crucial for accuracy and context and shows depth of engagement for transparency.

\paragraph{Agentic Systems:}
For agentic systems that autonomously plan and execute tasks, the developed plans \should be reported as \supplementarymaterial if available, as they document the sequence of actions the system chose to pursue and the intermediate goals it set.
Interaction logs \should be generalized to include full interaction traces between humans and agentic systems, capturing human-in-the-loop feedback, approval or rejection decisions, and iterative refinements.
These traces \should also be reported as \supplementarymaterial so that readers can reconstruct which actions the agent took, verify human oversight decisions, and assess the quality of tool-generated outputs.

Agentic coding tools can be configured using version-controlled Markdown files such as \texttt{AGENTS.md}, which contain persistent, project-specific instructions~\cite{mohsenimofidi2026context}.
These \emph{context files} are automatically injected into prompts and can significantly influence agent behavior by specifying coding conventions, architectural constraints, tool preferences, and task-specific instructions.
Researchers \should report all context files used to configure AI agent behavior as part of their \supplementarymaterial, as these files are a form of prompt artifact that shapes agent behavior similarly to system prompts.

\guidelinesubsubsection{Example(s)}

A paper by \citeauthor{anandayuvaraj2024fail}~\cite{anandayuvaraj2024fail} is a good example of making prompts available online.
In that paper, the authors analyze software failures reported in news articles and use prompting to automate tasks such as filtering relevant articles, merging reports, and extracting detailed failure information.
Their online appendix contains all the prompts used in the study, providing transparency and supporting reproducibility.

\citeauthor{Liang2024}'s paper is another good example of comprehensive prompt reporting.
The authors make the exact prompts available in their \supplementarymaterial on Figshare, including details such as code blocks being enclosed in triple backticks.
The \paper thoroughly explains the rationale behind the prompt design and the data output format.
It also includes an overview figure and two concrete examples, enabling transparency and reproducibility while keeping the main text concise.

An example of reporting full interaction logs is the study by \citeauthor{ronanki2023investigating}~\cite{ronanki2023investigating}, for which they reported the full answers of ChatGPT and uploaded them to \href{https://zenodo.org/records/8124936}{Zenodo}. 

\guidelinesubsubsection{Benefits}

Detailed prompt documentation improves verifiability, reproducibility, and comparability of LLM-based studies, allowing other researchers to replicate studies, refine prompts, and evaluate how different content types and formatting choices influence LLM behavior.

Unlike human participant conversations, which often cannot be reported due to confidentiality, LLM interaction logs can be shared. This enables reproduction studies, tracking of response changes over time or across model versions, and secondary research on LLM consistency for specific SE tasks.

\guidelinesubsubsection{Challenges}

One challenge is the complexity of prompts that combine multiple components, such as code, error messages, and explanatory text.
Formatting differences (e.g., Markdown vs. plain text) can affect how LLMs interpret input.
Additionally, prompt length constraints may require careful context management, particularly for tasks involving extensive artifacts such as large codebases.
Privacy and confidentiality concerns can hinder prompt sharing, especially when sensitive data is involved.

Not all systems allow the reporting of complete (system) prompts and interaction logs with ease.
This hinders transparency and verifiability.
We also recommend exploring open-source tools such as \emph{OpenCode}~\cite{opencode}.
For commercial tools, researchers \must report all available information and acknowledge unknown aspects as limitations.
Understanding suggestions of commercial tools such as \emph{GitHub Copilot} might require recreating the exact state of the codebase at the time the suggestion was made, a challenging context to report.
One solution could be to use version control to capture the exact state of the codebase when a recommendation was made, keeping track of the files that were automatically added as context.

\guidelinesubsubsection{Study Types}

Researchers \must report the complete prompts that were used, including all instructions, context, and input data. This applies across all study types, but reporting requirements and specific focus areas vary.

For \newtools, researchers \must explain how prompts were generated and structured within the tool.
For \llmusage, especially for controlled experiments, exact prompts \must be reported for all conditions.
For observational studies, especially the ones targeting commercial tools, researchers \must report the full interaction logs except for when transcripts might identify anonymous participants or reveal personal or confidential information.
As mentioned before, if the complete interaction logs cannot be shared, e.g., because they contain confidential information, the prompts and responses \must at least be summarized and described in the \paper.
For \annotators, researchers \must document any predefined coding guides or instructions included in prompts, as these influence how the model labels artifacts.
For \judges, researchers \must report the evaluation criteria, scales, and examples embedded in the prompt to ensure a consistent interpretation.
For \synthesis tasks (e.g., summarization, aggregation), researchers \must document the sequence of prompts used to generate and refine outputs, including follow-ups and clarification queries.
For \subjects (e.g., simulating human participants), researchers \must report any role-playing instructions, constraints, or personas used to guide LLM behavior.
For \benchmarkingtasks studies using pre-defined prompts (e.g., \emph{HumanEval}, \emph{SWE-Bench}), researchers \must specify the benchmark version and any modifications made to the prompts or evaluation setup.
If prompt tuning, retrieval-augmented generation (RAG), or other methods were used to adapt prompts, researchers \must disclose and justify those changes, and \should make the relevant code publicly available.

%% Suggestion: Consider integrating the footnote defining "over-rationalized" into the main text for smoother reading.
%% Suggestion: Consider splitting into two paragraphs: (1) what NOT to demand from authors; (2) what TO focus on (reproducibility + bias transparency).
\guidelinesubsubsection{Advice for Reviewers}

As with other guidelines, missing prompt information is typically a minor revision request unless so much is missing that methodological rigor cannot be assessed. The most challenging aspect is the description of prompt development, which is inherently iterative and creative. Reviewers should \textit{not} expect justification of each word choice or post-hoc rationalized accounts of prompt generation. Instead, reviewers should focus on whether (1) a new research team could \textit{use} (not reproduce) exactly the same prompts in exactly the same way; and (2) potential biases or validity issues are transparent.

\guidelinesubsubsection{See Also}
\begin{itemize}[label=$\rightarrow$]
    \item Section~\modelversion: model versions and configuration for prompt-model mapping.
    \item Section~\toolarchitecture: agent behavior traceability requirements.
    \item Section~\humanvalidation: evaluation of agentic tool outputs.
    \item Section~\openllm: open-source tools for transparent prompt logging.
    \item Section~\limitationsmitigations: reporting limitations of commercial tool transparency.
\end{itemize}
 
\subsection{Use Human Validation for LLM Outputs (G5)}
\label{sec:use-human-validation-for-llm-outputs}

\vspace{0.5\baselineskip}
\begin{framed}
\tldr: If assessing the quality of generated artifacts is important and no reference datasets or suitable comparison metrics exist, researchers \should use human validation for LLM outputs. If they do, they \must define the measured construct (e.g., usability, maintainability) and describe the measurement instrument in the \paper. Researchers \should consider human validation early in the study design (not as an afterthought) and build on established reference models for human-LLM comparison. When developing or adapting measurement instruments, researchers \must share them. When aggregating LLM judgments, methods and rationale \should be reported and inter-rater agreement \should be assessed. Confounding factors \should be controlled for, and power analysis \should be conducted to ensure statistical robustness. \end{framed}

\guidelinesubsubsection{Rationale}

Automated metrics alone cannot fully capture the validity of subjective or complex constructs.
When LLMs are used to automate tasks that involve human judgment, their outputs must be validated against human assessments to ensure that the results are not only reliable but also valid.

\guidelinesubsubsection{Recommendations}

When LLMs are used to automate research or software development tasks previously performed by humans, it is essential to assess the LLM's performance. When existing reference datasets or comparison metrics cannot fully capture the target qualities relevant in the study context, researchers \should rely on human judgment to validate the LLM outputs.

\paragraph{Study Design Considerations:}
Integrating human participants in the study design will require additional considerations, including a recruitment strategy, annotation guidelines, training sessions, or ethical approvals.
Therefore, researchers \should consider human validation early in the study design, not as an afterthought.
Authors \must clearly define in the \paper the constructs that the human and LLM annotators evaluate~\cite{DBLP:conf/ease/RalphT18}.
When designing custom instruments to assess LLM output (e.g., questionnaires, scales), researchers \must share their instruments.

\paragraph{Subjective Judgment and Agreement:}
When the judgment is subjective (i.e., depends on the judge's values or theories), the LLM output \should be evaluated against judgments aggregated from multiple human judges, and researchers \should clearly describe their aggregation method and reasoning. The preferred method is to randomly order the objects requiring judgment and put them in groups small enough to judge in one to three hours. Two or three human experts review the first group together, simultaneously judging, discussing, and creating a set of decision rules documenting their reasoning. Then, the experts iterate among rounds of:
\begin{enumerate}
    \item independently rating one group of objects,
    \item calculating inter-rater agreement (IRA) or reliability (IRR),
    \item meeting to discuss disagreements and reach consensus,
    \item updating the decision rules so the disagreement will not recur.
\end{enumerate}

The goal is to reach a target IRA or IRR (e.g., Krippendorff's $\alpha>0.8$), indicating sufficient decision rules. Once this target is reached and sustained for two to three groups, it is permissible to continue with a single rater.
Researchers \should report measures of IRA or IRR, preferably broken down by round.

Confounding factors \should be discussed and, where feasible, controlled for (e.g., by categorizing participants according to their level of experience or expertise).
Where applicable, researchers may perform a power analysis~\cite{Cohen1992,DBLP:journals/infsof/DybaKS06} to estimate the required sample size, ensuring sufficient statistical power in their experimental design.
However, we also note that to our knowledge, there is limited established guidance specific to determining sample sizes for LLM-human comparisons. In other fields, 100 comparisons are often made without further explanation~\cite{tam2024framework}.

Researchers \should use established reference models to compare humans with LLMs.
For example, \citeauthor{Schneider2025ReferenceModel}~\cite{Schneider2025ReferenceModel} outline design considerations for studies comparing LLMs with humans.
If studies aim to automate annotating software artifacts using an LLM, researchers \should follow systematic approaches to decide whether and how human annotators can be replaced (e.g., showing that a jury of three LLMs exhibit model-to-model agreement comparable to trained human annotators, such as Krippendorff's $\alpha>0.8$). Model-to-model agreement should not be held to a lower standard than human-to-human agreement, as low thresholds may simply reflect shared model biases rather than annotation quality.

\paragraph{Agentic Tools}
Besides generating and modifying content, agentic software development tools such as \emph{Claude Code} can autonomously call command-line tools or pull in additional information from MCP servers.
For such tools, it makes sense to separate textual output (e.g., file changes) from output related to tool execution.
When evaluating agentic tools, researchers \should assess the feedback that users provided for proposed changes, report statistics on how frequently they accepted the content, and how they modified it.
This agentic human-in-the-loop interaction approach is essentially a built-in human validation, even though the degrees of freedom are larger than in more traditional experiments that use LLMs directly.

\guidelinesubsubsection{Example(s)}

\citeauthor{DBLP:conf/msr/AhmedDTP25}~\cite{DBLP:conf/msr/AhmedDTP25} proposed a systematic method for deciding whether LLMs can replace human annotators on a given task, using model-to-model agreement as an initial screening criterion (with a threshold of $\alpha>0.5$) and model confidence for sample-level decisions. However, their threshold is well below the levels generally considered acceptable for inter-rater agreement.
\citeauthor{Krippendorff2018} recommends discarding data with $\alpha<0.667$, considers $0.667 \leq \alpha < 0.8$ sufficient only for tentative conclusions, and requires $\alpha \geq 0.8$ for reliable data~\cite{Krippendorff2018}. Researchers adopting similar screening approaches should apply thresholds consistent with these established standards. Notably, while three LLMs exhibited higher inter-model agreement than human annotators on some tasks, human-model agreement remained low on others.
This illustrates that high model-model reliability does not guarantee alignment with human judgment, reinforcing the need for human validation before assuming that LLM annotations are valid.
Moreover, a high model-model agreement could merely indicate that the models share systematic biases and hence reliably agree on the wrong answer.

%Meanwhile, \citeauthor{DBLP:conf/icse/XueCBTH24}~\cite{DBLP:conf/icse/XueCBTH24} conducted a controlled experiment to evaluate the impact of ChatGPT on the performance and perceptions of students in an introductory programming course.
%They used multiple measures to judge LLM impacts from the human point of view including recording students' screens, evaluating their answers for given tasks, and distributing an opinion survey.
\citet{hymel2025analysisllmsvshuman}~\cite{hymel2025analysisllmsvshuman} evaluated ChatGPT's capability to generate requirements documents by comparing an LLM-generated and a human-generated document based on the same business use case. Domain experts reviewed both documents and attempted to distinguish their origin.% and judged in terms of alignment with the original business use case (scale of 1-10), requirements completion (Not Complete, Fairly Complete, Fully Complete), and whether they believed it was created by a human or an LLM.
%Finally, they analyzed the influence of the participants' familiarity with AI tools on the study results.
%Participants self-reported their AI tool proficiency (Novice, Intermediate, Advanced, Expert) and usage frequency (Daily, Weekly, Monthly, Never), which were then correlated with their ratings of the requirements documents.

\guidelinesubsubsection{Benefits}

Validating LLMs against human judgments builds confidence in the LLM's accuracy and validity, increasing the trustworthiness of findings, especially when IRA or IRR metrics are reported~\cite{khraisha2024canlargelanguagemodelshumans}. Furthermore, incorporating feedback from human judges may help to improve the LLM or LLM-based tool based on the reported experiences.

\guidelinesubsubsection{Challenges}

Assessing a construct like LLM performance is challenging because ensuring that a construct is defined well and operationalized using an appropriate measurement model requires a deep understanding of (1) the construct, (2) construct validity in general, and (3) instrumentation~\cite{DBLP:journals/tse/SjobergB23,ralph2024teaching}. Comparing an LLM to human judges is typically slower and more expensive than machine-generated measures. More fundamentally, neither human judgment nor machine-generated measures provides an objective ground truth against which LLM accuracy can be firmly determined.   

Human judgments exhibit variability due to differences in experience, expertise, interpretations, and personal biases~\cite{DBLP:journals/pacmhci/McDonaldSF19}. When diverse humans rate items reliably given clear decision rules, we assume that reliability implies validity, but it does not. Measuring reliability is \textit{much} easier than measuring validity, and often the best researchers can do is argue conceptually for why their judges, decision rules, and constructs \textit{should} produce valid ratings.

Researchers must weigh the cost and time need for human validation against the expected corresponding improvement in validity over machine-generated measures. 

%For example, LLM results vary significantly in natural language processing tasks, casting doubt on their reliability and the possibiltiy of replacing human judges~\cite{DBLP:journals/corr/abs-2406-18403}. Moreover, LLMs do not match human annotation in labeling tasks for natural language inference, position detection, semantic change, and hate speech detection~\cite{DBLP:conf/chi/Wang0RMM24}.

\guidelinesubsubsection{Study Types}

This guideline applies to all study types, although the need for human validation varies.
For \annotators, researchers \should validate LLM-generated annotations against human annotators to assess labeling quality and identify systematic biases.
When using \judges, researchers \should co-create initial rating criteria with humans and validate a sample of LLM judgments against human expert assessments.
For \synthesis, researchers \should employ human oversight to verify that qualitative interpretations and synthesized outputs faithfully represent the underlying data.
For \subjects, researchers \should validate simulated responses against real human data to assess the fidelity of the simulation.
For \llmusage, researchers \should carefully reflect on the validity of their evaluation criteria and validate subjective assessments with human experts.
In developing \newtools, human validation of the LLM output is critical when the output is supposed to match human expectations.
For \benchmarkingtasks, there is less need for human validation when using extensively validated and widely-used benchmarks, but researchers \should employ human validation when creating or adapting new benchmarks.

\guidelinesubsubsection{Advice for Reviewers}

Human validation may be the most challenging of our guidelines to assess because it often requires evaluating conceptual arguments. If LLM output is validated only by comparison with other LLMs, reviewers should look for \textit{quantitative empirical evidence} that such comparison is reliable \textit{and} valid. High inter-model agreement alone is insufficient, as reliability does not imply validity. Similarly, reviewers should expect evidence that any employed benchmarks are reliable and valid. Absent such evidence, human validation is warranted.
A single human judge is appropriate only when judgments depend on widely accepted theories and involve limited value conflict (e.g., tagging method names containing abbreviations). For multiple judges, reviewers should expect IRA/IRR improvement techniques as described in the recommendations above (experienced raters, organized rounds, consensus meetings, updated decision rules). Low IRA or IRR (e.g., Krippendorff's $\alpha<0.8$) without these techniques is a concern. Conversely, if authors have followed best practices and still obtained mediocre results (e.g., $0.66<\alpha<0.8$), this should be noted as a limitation.
Beyond reliability, reviewers should expect authors to explain conceptually why their human judgments should be valid, considering construct definitions, decision rules, and judge expertise.
As with other guidelines, missing information is typically a revision request. Absent judge instructions, instruments, decision rules, or construct definitions may prevent assessment of rigor and validity, while missing recruitment details are less critical. Clarification requests about construct definitions are routine and should not alone warrant rejection. 
\subsection{Use an Open LLM as a Baseline (G6)}
\label{sec:use-an-open-llm-as-a-baseline}

\vspace{0.5\baselineskip}
\begin{framed}
\tldr: Researchers \should include an open LLM as a baseline when using commercial models and report inter-model agreement. By open LLM, we mean a model that everyone with the required hardware can deploy and operate. Such models are usually ``open weight.'' A full replication package with step-by-step instructions \should be provided as part of the \supplementarymaterial. \end{framed}

\guidelinesubsubsection{Rationale}

Reproducibility depends on access to the model under study.
When research relies exclusively on proprietary models, other researchers cannot independently verify or build upon the findings.
Including an open LLM as a baseline ensures that at least part of the study can be fully replicated.

\guidelinesubsubsection{Recommendations}

Empirical studies using LLMs in SE, especially those that target commercial tools or models, \should incorporate an open LLM as a baseline and report established metrics for inter-model agreement.
We acknowledge that including an open LLM baseline might not always be possible, for example, if the study involves human participants, and letting them work on the tasks using two different models might not be feasible.
Using an open model as a baseline is also not necessary if the use of the LLM is tangential to the study goal. %, or strict reproducibility of the LLM portion of the results is not necessary.

Open models allow other researchers to verify research results and build upon them, even without access to commercial models.
A comparison of commercial and open models also allows researchers to contextualize model performance.
Researchers \should provide a complete replication package as part of their \supplementarymaterial, including clear step-by-step instructions on how to verify and reproduce the results reported in the paper.

Open LLMs are available on platforms such as \href{https://huggingface.co/}{\emph{Hugging Face}} and can be hosted locally using frameworks such as \href{https://ollama.com/}{\emph{Ollama}} or \href{https://lmstudio.ai/}{\emph{LM Studio}}, or on cloud services such as \href{https://together.ai/}{\emph{Together AI}}, AWS, Azure, and Google Cloud.

The term ``open'' can have different meanings in the context of LLMs.
\citeauthor{widder2024open} discuss three types of openness (transparency, reusability, and extensibility) and what openness can and cannot provide~\cite{widder2024open}.
The \emph{Open Source Initiative} (OSI)~\cite{OSIAI2024} defines open-source AI as having access to everything needed to understand, modify, share, retrain, and recreate the model.

Researchers can also explore open-source tools such as Continue, Cline, and opencode as alternatives to commercial tools like GitHub Copilot and Claude Code, enabling instrumentation, architectural transparency, and detailed telemetry collection.

\guidelinesubsubsection{Example(s)}

An increasing number of studies have adopted open LLMs as baseline models.
For example, \citeauthor{wang2024and} evaluated seven advanced LLMs, six of which were open-source, testing 145 API mappings drawn from eight popular Python libraries across 28,125 completion prompts aimed at detecting deprecated API usage in code completion~\cite{wang2024and}.
\citeauthor{moumoula2024large} compared four LLMs on a cross-language code clone detection task~\cite{moumoula2024large}.
Three evaluated models were open-source.
\citeauthor{gonccalves2025evaluating} fine-tuned the open LLM \emph{LLaMA} 3.2 on a refined version of the \emph{DiverseVul} dataset to benchmark vulnerability detection performance~\cite{gonccalves2025evaluating}.
\emph{CodeBERT}, a bimodal transformer pre-trained by Microsoft Research, has been widely used as an open baseline~\cite{DBLP:journals/jss/YangZCZHC23, DBLP:conf/gaiis/XiaSD24, DBLP:conf/kbse/SonnekalbGBM22, DBLP:conf/icse/CaiYMMN24}, with model weights, source code, and data-processing scripts published on GitHub~\cite{codebert}.

\guidelinesubsubsection{Benefits}

Using a true open LLM as a baseline improves the reproducibility of scientific research by providing access to model architectures and parameter settings, and ideally training data, thereby allowing independent reconstruction and verification of experimental results.
Moreover, by adopting an open-source baseline, researchers can directly compare novel methods against established performance metrics without the variability introduced by proprietary systems.
The transparent nature of these models allows for detailed inspection of data processing pipelines and decision-making routines, which is essential for identifying potential sources of bias and delineating model limitations.
Furthermore, unlike closed-source alternatives, which can be withdrawn or altered without notice, open LLMs ensure long-term accessibility and stability, preserving critical resources for future studies.
Finally, the licensing requirements associated with open-source implementations lower financial barriers, making advanced language models attainable for research groups operating under constrained budgets.

\guidelinesubsubsection{Challenges}

Open-source LLMs face several notable challenges.
First, they often lag behind the most advanced proprietary ``frontier'' models in common benchmarks, making it difficult for researchers to demonstrate clear improvements when evaluating new methods using open LLMs alone.
Additionally, deploying and experimenting with these models typically requires substantial hardware resources, in particular high-performance GPUs, which may be beyond reach for many academic groups.
The notion of ``openness'' remains in flux: many models release only trained weights without training data or methodological details (``open weight'' openness)~\cite{Gibney2024}, which is why we reference the OSI definition in our recommendations~\cite{OSIAI2024}.
Finally, unlike APIs provided by proprietary vendors (e.g., the OpenAI API), installing, configuring, and fine-tuning open-source models can be technically demanding.
Documentation is often sparse or fragmented, placing a high barrier to entry for researchers without specialized engineering support.

\guidelinesubsubsection{Study Types}

This guideline applies primarily to study types in which the researcher controls which LLM is used.
In formal benchmarking studies and controlled experiments (see \benchmarkingtasks), an open LLM \must be one of the models under evaluation.
When evaluating \newtools, researchers \should use an open LLM as a baseline whenever it is technically feasible; if integration proves too complex, they \should report the initial benchmarking results of open models.
For \annotators and \judges, researchers \should compare annotation or judgment quality from open vs.\ commercial models to assess the extent to which results depend on a specific proprietary model.
For \synthesis, researchers \should compare synthesis results from open and commercial models to evaluate robustness of the findings.
For \llmusage, using an open LLM as a baseline is often not feasible when the study observes participants using specific commercial tools; in such cases, investigators \should explicitly acknowledge its absence and discuss how this limitation might affect their conclusions.
For \subjects, using an open LLM as a baseline may similarly be impractical if the study design requires the capabilities of a specific model; researchers \should acknowledge this limitation when applicable.

\guidelinesubsubsection{Advice for Reviewers}

Reviewers should distinguish between LLM use that is central to the research (e.g., building LLM-driven SE tools, using LLMs for synthesis) and use that is tangential (e.g., generating recruiting materials). An open LLM baseline is expected only when LLM use is central; otherwise, its absence need not be justified.
When an open baseline is expected, reviewers should look for either the use of an open LLM or a convincing argument for why it is impractical. If authors claim openness, some justification for that characterization is appropriate. Where practical, reviewers should examine whether the replication package contains sufficiently detailed instructions.
Reviewers should \textit{not} penalize studies for performance differences between open and proprietary models, as this is beyond the authors' control.

\guidelinesubsubsection{See Also}
\begin{itemize}[label=$\rightarrow$]
    \item Section~\modelversion: model version, configuration, and parameter settings.
    \item Section~\toolarchitecture: hardware and hosting requirements for open models.
    \item Section~\benchmarksmetrics: metrics for inter-model agreement and benchmarking.
    \item Section~\limitationsmitigations: reporting limitations when open baselines are absent.
\end{itemize}
 
\subsection{Use Suitable Baselines, Benchmarks, and Metrics (G7)}
\label{sec:use-suitable-baselines-benchmarks-and-metrics}

\vspace{0.5\baselineskip}
\begin{framed}
\tldr: Researchers \must justify all benchmark and metric choices in the \paper and \should summarize benchmark structure, task types, and limitations. Where possible, traditional (non-LLM) baselines \should be used for comparison. Researchers \must explain in the \paper why the selected metrics are suitable for the specific study; prior adoption in related work alone does not constitute sufficient justification. They \should report established metrics to make study results comparable, but can report additional metrics that they consider appropriate. Due to the inherent non-determinism of LLMs, experiments \should be repeated; the result distribution \should then be reported using descriptive statistics. If comparing models or tools, researchers \must use appropriate inferential statistics rather than relying solely on summary statistics. Researchers \should justify the number of experiment repetitions, for example through a power analysis or by monitoring convergence of descriptive statistics. \end{framed}

% \todo{Paul Ralph: Moving this guideline (use suitable baselines, etc.) above guideline 5 (Use human validation) would enable a cleaner presentation because we'd be able to explain about validating benchmarks here, then say "and if you don't have a validated benchmark then use humans".}

\guidelinesubsubsection{Rationale}

Meaningful evaluation requires well-understood, valid measurement instruments.
Without justified benchmarks and metrics, claims about LLM performance lack the rigor needed for scientific comparison and cumulative progress.
Benchmarks, baselines, and metrics play an important role in assessing the effectiveness of LLMs and LLM-based tools. A \textit{benchmark} is a ``standard tool for the competitive evaluation and comparison of competing systems or components according to specific characteristics, such as performance, dependability, or security''~\cite{Kistowski2015Benchmark}. Meanwhile, ``a \textit{metric} is a method, algorithm, or procedure for assigning one or more numbers to a phenomenon''~\cite{ralph2024teaching}. A baseline is a reference point, enabling comparison of LLMs against traditional algorithms with lower computational costs.

\guidelinesubsubsection{Recommendations}

When selecting benchmarks and metrics, it is important to fully understand the benchmark tasks, what exactly is being measured, and how it relates to the (often latent) variables researchers actually care about. If one or more metrics or benchmarks are used, researchers:
\begin{itemize}
    \item \must briefly justify (in the \paper) why selected benchmarks and metrics are suitable for the given task or study;
    \item \must discuss the reliability and validity (especially construct validity) of selected metrics and benchmarks;
    \item \should summarize the structure and tasks of the selected benchmark(s), including the programming language(s) and descriptive statistics such as the number of contained tasks and test cases;
    \item \should discuss the limitations of the selected benchmark(s) (e.g. widely used benchmarks such as \emph{HumanEval} \cite{DBLP:journals/corr/abs-2107-03374} and \emph{MBPP} \cite{DBLP:journals/corr/abs-2108-07732} Python functions assess a very specific part of software development, which is not representative of the full breadth of SE work~\cite{Chandra2025benchmarks}).
    \item \should include an example of a task and the corresponding test case(s) to illustrate the structure of the benchmark.
\end{itemize}

If multiple benchmarks exist for the same task, researchers \should compare performance between benchmarks.
When selecting only a subset of all available benchmarks, researchers \should use the most specific benchmarks given the context.

Furthermore, researchers \should check whether a less resource-intensive approach (e.g., for static analysis tasks or program repair) can serve as a baseline. If so, the LLM or LLM-based tool should be compared with such baselines using suitable metrics. Even if LLM-based tools outperform baselines, researchers \should discuss whether the resources consumed justify the (often marginal) improvements~\cite{DBLP:journals/cacm/Menzies25}.

To compare traditional and LLM-based approaches or different LLM-based tools, researchers \should report established metrics whenever possible, as this allows secondary research.
They can, of course, report additional metrics that they consider appropriate.
We briefly discuss common metrics in the Example(s) subsection below.
As mentioned, researchers \must motivate why they chose a certain metric or variant thereof for their particular study.
Prior adoption alone does not constitute sufficient justification; researchers \mustnot justify metric choices solely by citing their use in prior work.

Due to LLM non-determinism, researchers \should repeat experiments and report descriptive statistics of model or tool performance (e.g., arithmetic mean, median, confidence intervals, standard deviations)~\cite{DBLP:conf/nips/AgarwalSCCB21, bjarnason2026randomnessagenticevals}. If comparing models or tools, researchers \must use appropriate inferential statistics (e.g., hypothesis tests such as the Mann-Whitney U test, McNemar's test for binary outcomes~\cite{DBLP:conf/iclr/KublerBKZYKK26}, or bootstrap-based comparisons, along with effect sizes) rather than relying solely on differences in means or other summary statistics.

The number of required repetitions depends on factors such as the study type, the variability of the task, and the desired precision of estimates. As with sample sizes for human validation (\humanvalidation), researchers \should justify their chosen number of repetitions, for example, through a power analysis~\cite{Cohen1992, bjarnason2026randomnessagenticevals} or by monitoring the convergence of descriptive statistics across incremental runs~\cite{blackwell2024towards}. A pilot study can help estimate the expected variability and inform this decision.

From a measurement perspective, researchers \should reflect on the theories, values, and measurement models on which the benchmarks and metrics they have selected for their study are based.
For example, the phenomena labeled ``bugs'' in a large open dataset of software bugs are according to a certain theory of what constitutes a bug, and from the values and perspectives of the people who labeled the dataset. Reflecting on the context in which these labels were assigned and discussing whether and how the labels generalize to a new study context is crucial.

\guidelinesubsubsection{Example(s)}

\paragraph{Common Metrics:}

Two common metrics used for generation tasks are \emph{BLEU-N} and \emph{pass@k}.
\emph{BLEU-N}~\cite{DBLP:conf/acl/PapineniRWZ02} was originally developed for evaluating machine translation quality by measuring n-gram precision between a candidate and reference text, ranging from $0$ (dissimilar) to $1$ (similar). It has been widely adopted in SE for code generation tasks, though its validity in this context is debatable: n-gram overlap does not capture functional correctness, and syntactically different code can be semantically equivalent (see Challenges below).
Code-specific variations such as \emph{CodeBLEU}~\cite{DBLP:journals/corr/abs-2009-10297} and \emph{CrystalBLEU}~\cite{DBLP:conf/kbse/EghbaliP22} attempt to address these limitations by introducing additional heuristics such as AST matching.

The metric \emph{pass@k} reports the likelihood that a model correctly completes a code snippet at least once within $k$ attempts. Originally introduced as \emph{success rate at B} by \citeauthor{DBLP:journals/corr/abs-1906-04908}~\cite{DBLP:journals/corr/abs-1906-04908} and popularized by \citeauthor{DBLP:journals/corr/abs-2107-03374}~\cite{DBLP:journals/corr/abs-2107-03374}, it ranges from $0$ (no correct solution in $k$ tries) to $1$ (at least one correct solution), with correctness typically defined by passing test cases.

The \emph{pass@k} metric is defined as:
\[
\text{pass@}k = 1 - \frac{\binom{n-c}{k}}{\binom{n}{k}}\,, \text{where:}
\]
\begin{itemize}
  \item $n$ is the total number of generated samples per prompt,
  \item $c$ is the number of correct samples among $n$, and
  \item $k$ is the number of samples drawn.
\end{itemize}

The choice of $k$ depends on the downstream task: \emph{pass@1} is critical for single-suggestion scenarios like code completion, while higher $k$ values (e.g., $2$, $5$, $10$) assess multi-attempt capability~\cite{DBLP:journals/corr/abs-2308-12950, DBLP:journals/corr/abs-2401-14196, DBLP:journals/corr/abs-2409-12186, DBLP:journals/corr/abs-2305-06161}.
However, \emph{pass@k} is not a universal metric suitable for all generation tasks.
It requires a binary notion of correctness, making the metric appropriate for code synthesis evaluated via unit tests, but unsuitable for open-ended generation tasks such as comment generation, where multiple valid outputs exist.

If a study evaluates an LLM-based tool for supporting humans, a relevant metric is the acceptance rate, meaning the ratio of all accepted artifacts (e.g., test cases, code snippets) in relation to all artifacts that were generated and presented to the user.
Another way of evaluating LLM-based tools is calculating inter-model agreement.
This allows researchers to assess how dependent a tool's performance is on specific models and versions.
Metrics used to measure inter-model agreements include general agreement (percentage), \emph{Cohen's kappa}, and \emph{Krippendorff's $\alpha$} (see \humanvalidation for recommended thresholds and best practices for measuring agreement).

Common problem types for LLM-based studies are classification, recommendation, and generation, each requiring different metrics~\cite{10.1145/3695988}. \citeauthor{hu2025assessingadvancingbenchmarksevaluating}~\cite{hu2025assessingadvancingbenchmarksevaluating} categorized 191 LLM benchmarks by SE task, providing a valuable reference. Common metrics include \emph{BLEU}, \emph{pass@k}, \emph{Accuracy@k}, and \emph{Exact Match} for generation; \emph{Mean Reciprocal Rank} for recommendation; and \emph{Precision}, \emph{Recall}, \emph{F1-score}, and \emph{Accuracy} for classification.

\paragraph{Benchmark Examples:}

Benchmarks used for code generation include \emph{HumanEval} (available on \href{https://github.com/openai/human-eval}{GitHub}) \cite{DBLP:journals/corr/abs-2107-03374}, \emph{MBPP} (available on \href{https://huggingface.co/datasets/google-research-datasets/mbpp}{Hugging Face}) \cite{DBLP:journals/corr/abs-2108-07732},
% Both benchmarks consist of code snippets written in Python sourced from publicly available repositories.
% Each snippet consists of four parts: a prompt containing a function definition and a corresponding description of what the function should accomplish, a canonical solution, an entry point for execution, and test cases.
% The input of the LLM is the entire prompt.
% The output of the LLM is then evaluated against the canonical solution using metrics or against a test suite.
\emph{ClassEval} (available on \href{https://github.com/FudanSELab/ClassEval}{GitHub}) \cite{DBLP:journals/corr/abs-2308-01861}, \emph{LiveCodeBench} (available on \href{https://github.com/LiveCodeBench/LiveCodeBench}{GitHub}) \cite{DBLP:journals/corr/abs-2403-07974}, and \emph{SWE-bench} (available on \href{https://github.com/swe-bench/SWE-bench}{GitHub}) \cite{DBLP:conf/iclr/JimenezYWYPPN24}.
An example of a code translation benchmark is \emph{TransCoder}~\cite{DBLP:journals/corr/abs-2006-03511} (available on \href{https://github.com/facebookresearch/CodeGen}{GitHub}). 

Most benchmarks focus on generation tasks, but benchmarks for classification and recommendation also exist.
For classification, \emph{DiverseVul} (available on \href{https://github.com/wagner-group/diversevul}{GitHub})~\cite{DBLP:conf/raid/0001DACW23} provides vulnerable and non-vulnerable functions evaluated using standard classification metrics.
For recommendation, \emph{CodeSearchNet} (available on \href{https://github.com/github/CodeSearchNet}{GitHub})~\cite{DBLP:journals/corr/abs-1909-09436} contains code-documentation pairs evaluated using \emph{Mean Reciprocal Rank}.

\guidelinesubsubsection{Benefits}

Benchmarks and metrics improve reproducibility and comparability across studies, leading to faster improvements in the cumulative body of knowledge. It is simply easier to assess a new system against one or a few benchmarks than hundreds of competing systems, and when most systems are assessed against the same benchmarks and metrics, we can see which approaches work best (at least, on those benchmarks). This comparison enables progress tracking; for example, when researchers iteratively improve a new LLM-based tool and test it against benchmarks after significant changes. For practitioners, leaderboards (published benchmark results of models) support selecting models for downstream tasks. Meanwhile, baselines help us understand just how much LLMs improve performance over non-LLM-enabled alternatives. 

\guidelinesubsubsection{Challenges}

Computer scientists have a long history of using oversimplified, unvalidated metrics~\cite{DBLP:conf/ease/RalphT18} (e.g., lines of code, CPU time) as proxies for complex, multidimensional latent variables (e.g., system size, environmental sustainability). In fields such as psychology, where measurement theory has a longer tradition~\cite{borsboom2005measuring}, metrics are backed by foundational theories or extensive empirical validation of their psychometric properties. The fundamental challenge with AI metrics and benchmarks is that many have neither firmly understood theoretical underpinnings nor extensive empirical validation of their construct, measurement, and ecological validity.

For example, \emph{pass@k} is intended to measure a model's \textit{functional correctness}, that is, its ability to generate code that produces the expected output for a given specification. However, functional correctness is only one dimension of code quality. \emph{pass@k} does not capture maintainability, readability, security, or efficiency of the generated code, all of which are critical for downstream use. Furthermore, ``correctness'' is defined entirely by test suites, whose coverage is itself unvalidated: a solution that passes all provided tests may still be incorrect for untested inputs. Whether a test suite adequately operationalizes correctness for a given task is a subjective judgment that is rarely examined.

Similarly, \emph{HumanEval} is intended to measure a model's ability to \textit{synthesize short Python functions from docstring specifications}. This operationalizes a narrow slice of ``code generation ability'': it covers neither multi-file tasks, nor debugging, nor the use of existing codebases---tasks that dominate real-world software engineering~\cite{Chandra2025benchmarks}. Notably, neither \emph{pass@k} nor \emph{HumanEval} has undergone rigorous empirical validation of its construct validity; their widespread adoption rests on face validity and convenience rather than on evidence that they reliably measure what researchers intend them to measure~\cite{cao2025should}.

Benchmarks and metrics also have generalizability problems. Prominent LLM benchmarks such as \emph{HumanEval} and \emph{MBPP} use Python, so researchers can optimize for Python's idiosyncrasies, creating illusory improvements that do not generalize to other languages. Similarly, the metric \emph{BLEU-N} is a syntactic metric, so code can score highly without being executable. The metric \emph{Exact Match}, meanwhile, does not account for functional equivalence of syntactically different code. Both \emph{BLEU-N} and \emph{Exact Match} are influenced by code formatting, which confounds their intended use.

Execution-based metrics such as \emph{pass@k} directly evaluate correctness by running test cases, but they require a setup with an execution environment.
When researchers observe unexpected values for certain metrics, the specific results should be investigated in more detail to uncover potential problems.

% Many closed-source models, such as those released by \emph{OpenAI}, achieve exceptional performance on certain tasks but lack transparency and reproducibility~\cite{DBLP:conf/nips/00110ZZDJLHL24, DBLP:journals/corr/abs-2308-01861, DBLP:journals/corr/abs-2406-15877}.
% Benchmark leaderboards, particularly for code generation, are led by closed-source models~\cite{DBLP:journals/corr/abs-2308-01861, DBLP:journals/corr/abs-2406-15877}.
% While researchers \should compare performance against these models, they must consider that providers might discontinue them or apply undisclosed pre- or post-processing beyond the researchers' control (see \openllm).
%The lack of transparency introduces challenges in analyzing mistakes, limiting the understanding of model failures.
%Open-source models offer greater transparency since the entire process is under the researcher's control; however, they require appropriate hardware.

Finally, benchmark data contamination---where the benchmark itself is part of the training data---may lead to artificially high performance if the model remembers the solution from the training data rather than actually solving the task based on the input data~\cite{DBLP:journals/corr/abs-2406-04244}. Therefore, for proprietary LLMs that do not release their training data, researchers \should consider using human validation, curating new data, or refactoring existing data~\cite{DBLP:journals/corr/abs-2406-04244}. 

To mitigate contamination, researchers can create new benchmark datasets by collecting data after a specified cutoff date; researchers \must disclose data sources and collection dates for each release. However, this temporal approach requires continuous updates as new models may include the benchmark in future training data. Alternatively, keeping the benchmark private prevents inclusion in training sets, but requires trust in the benchmark creator and a system to execute it without leaking data. Researchers can also evaluate how contaminated their benchmark is~\cite{choi2025contaminatedbenchmarkquantifyingdataset}.


\guidelinesubsubsection{Study Types}

This guideline \must be followed for all study types that automatically evaluate the performance of LLMs or LLM-based tools.
The design of a benchmark and the selection of appropriate metrics are highly dependent on the specific study type and research goal.
Recommending specific metrics for specific study types is beyond the scope of these guidelines, but \citeauthor{hu2025assessingadvancingbenchmarksevaluating} provide a good overview of existing metrics for evaluating LLMs~\cite{hu2025assessingadvancingbenchmarksevaluating}.

For \benchmarkingtasks, this guideline is of primary importance: researchers \must use established benchmarks or rigorously justify the creation of new ones, and \must report standard metrics to enable cross-study comparison.
For \annotators, the research goal might be to assess which model comes close to a ground truth dataset created by human annotators.
Especially for open annotation tasks, selecting suitable metrics to compare LLM-generated and human-generated labels is important.
In general, annotation tasks can vary significantly.
Are multiple labels allowed for the same sequence? Are the available labels predefined, or should the LLM generate a set of labels independently?
Due to this task dependence, researchers \must justify their metric choice, explaining what aspects of the task it captures together with known limitations.
For \newtools, researchers \should benchmark the tool against suitable baselines using appropriate metrics that capture the tool's intended contribution.
For \judges, researchers \should report inter-rater agreement metrics and validity measures to demonstrate the reliability and quality of LLM judgments.
For \synthesis, researchers \should specify metrics for comparing synthesized outputs (e.g., coverage, faithfulness) and justify their appropriateness for the synthesis task.
For \llmusage, researchers \should justify the measurement instruments and metrics used for studying LLM usage patterns, including any survey scales or behavioral measures.
For \subjects, researchers \should compare simulated and real human responses using appropriate metrics to assess simulation fidelity.
If researchers assess a well-established task such as code generation, they \should report standard metrics such as \emph{pass@k} and compare the performance between models.
If non-standard metrics are used, researchers \must state their reasoning.

%% Changed: Softened "must not tolerate" → "should not accept" (advisory tone).
%% Changed: Simplified "ambivalence toward construct validity or obfuscation of validity concerns" → "dismissive or evasive treatment of construct validity".
%% Suggestion: Consider breaking the closing paragraph into shorter sentences for readability.
\guidelinesubsubsection{Advice for Reviewers}

Reviewers should expect manuscripts to:
\begin{enumerate*}
    \item clearly identify the constructs or variables the study aims to measure (e.g., LLM performance, quality of generated code), including independent, dependent, and control variables;
    \item present their measurement model, i.e., which metrics, benchmarks, or baselines are used and how they relate to the target constructs;
    \item justify the selection of any metrics, benchmarks, and baselines used;
    \item discuss \textit{in detail} the assumptions, reliability, and validity---\textit{especially construct validity}---of each benchmark and metric;
    \item articulate any limitations regarding construct and measurement validity.
\end{enumerate*}
As with other guidelines, missing information about baselines or metrics is typically a revision request. However, vague descriptions that conflate broad concerns (e.g., effectiveness, quality) with specific counting methods should be questioned. Ubiquity of a benchmark or metric does not imply validity or appropriateness for a given context. Manuscripts should convey a solid understanding of construct and measurement validity by explaining and justifying their measurement models.

\guidelinesubsubsection{See Also}
\begin{itemize}[label=$\rightarrow$]
    \item Section~\openllm: open models for inter-model comparison.
    \item Section~\humanvalidation: human evaluation to complement automated metrics.
    \item Section~\limitationsmitigations: reporting construct and measurement validity limitations.
\end{itemize} 
\subsection{Report Limitations and Mitigations (G8)}
\label{sec:report-limitations-and-mitigations}

\vspace{0.5\baselineskip}
\begin{framed}
\tldr: Researchers \must transparently report study limitations, including the impact of non-determinism and generalizability constraints. The \paper~\must specify whether generalization across LLMs or across time was assessed, and discuss model and version differences. Authors \must describe measurement constructs and methods, disclose any data leakage risks, and avoid leaking evaluation data into LLM improvement pipelines. They \must provide model outputs and discuss sensitive data handling. They \should justify LLM usage in light of its resource demands. Mitigation strategies such as replication packages, human validation, longitudinal re-runs, triangulation, and sensitivity analysis \should be employed and reported where applicable. Where full data sharing is not possible, a subset of the validation data \should be included to enable partial replication. \end{framed}

\guidelinesubsubsection{Rationale}

When using LLMs for empirical studies in SE, researchers face unique challenges and potential limitations that can influence the validity, reliability, and reproducibility of their findings~\cite{sallou2024breaking}.
It is important to openly discuss these limitations and explain how their impact was mitigated.
These limitations are relative to current LLM capabilities and tool architectures; speculating about future improvements is beyond the scope of a paper's limitation section.
Nevertheless, risk management and threat mitigation should be planned during study design, not as an afterthought.

\guidelinesubsubsection{Recommendations}
Researchers \must make every effort to present the limitations of their work clearly without defensiveness or obfuscation. These limitations may concern diverse topics including generalizability, internal validity (e.g. data leakage), reliability (i.e. non-determinism), and reproducibility (e.g. resource requirements). 
When deterministic reproducibility is structurally unattainable (e.g., SaaS-based models with opaque versioning), researchers \should adopt trustworthiness criteria from qualitative research to substantiate dependability and confirmability of findings.
LLM-based studies may involve many kinds of generalization including the following:
\begin{enumerate}
    \item From tested LLMs to others; that is, do the performance characteristics of the LLM(s) studied generalize to LLM(s) not included in the study.
    \item From tested configurations to others; that is, how sensitive are the results to the specific configuration(s) of the LLM under test. 
    \item From research to practice; for instance, just because researchers can get a certain level of code generation performance under ideal conditions does not mean software developers, given the same tools, will get similar results without having the researchers present to show them exactly what to do. 
    \item From a sample of people (e.g. human validators; human participants) to a larger population of people.
    \item From one time period to another; i.e., does the same LLM, or other versions of the same LLM, produce similar results at a different point in time?
\end{enumerate}

The following concerns present an overview that has to be tailored to the individual study context. This section does not repeat requirements from other recommendations.

\paragraph{Internal Validity:}
The primary threats to internal validity are:
\begin{itemize}
	\item Data leakage and contamination, including inter-dataset duplication, potentially resulting in a training–evaluation overlaps, leading to overly optimistic evaluation results.
	\item Unintended inclusion of evaluation data in model-improvement pipelines, contributing to data leakage. This is especially relevant for longitudinal studies including LLMs.
	\item Incomplete architecture, prompt, or pipeline reporting, posing hidden confounding factors.
\end{itemize}

Inter-dataset duplication is prevalent in SE, particularly for code-related benchmarks. As transparency on training data is limited for LLMs, researchers \must discuss potential data leakage effects and their impact on results.

\paragraph{Construct Validity:}
The primary threats to construct validity are:
\begin{itemize}
	\item Metric–construct mismatch (e.g., BLEU/ROUGE vs. functional correctness). More traditional metrics might not capture all relevant SE-specific aspects.
	\item Over-reliance on benchmark-specific metrics. Optimizing for a single benchmark might result in dataset-specific heuristics rather than the intended construct, overstating real-world utility.
	\item Benchmark scope limitations. Benchmarks commonly ignore runtime behaviors, security implications, readability, testability, and maintainability, yielding results that may not transfer to realistic development settings.
\end{itemize}

If constructs are based on subjective interpretations, purely automated metrics are insufficient. Researchers \must discuss how they ensured quality of subjective results, similarly to qualitative research.

\paragraph{External Validity:}
The primary threats to external validity are:
\begin{itemize}
    \item Cross-model transfer limitations. Results obtained with one LLM (or family of LLMs) may not generalize to others due to differences in training data, architecture, and capabilities.
    \item Tool-architecture specificity. Tools built around a specific LLM's API or capabilities may not transfer to other models without significant re-engineering.
    \item Limited domain coverage. Studies often focus on a narrow set of programming languages, task types, or application domains, limiting generalizability to other SE contexts.
    \item Limited participant diversity. Study participants (e.g., human validators, developers) may not represent the broader population in terms of expertise, geographic location, or cultural background.
    \item Cross-time instability (model evolution). Performance of proprietary models can change over time, leading to non-generalizable study outcomes~\cite{DBLP:journals/corr/abs-2307-09009, doi:10.1148/radiol.232411}.
\end{itemize}

Generalizability is particularly critical for proprietary and non-deterministic systems whose behavior is subject to drift (i.e., silent changes in model output over time)~\cite{DBLP:journals/corr/abs-2307-09009}. The limitations and mitigations used for external validity \must be discussed by the researchers.

\paragraph{Reliability \& Reproducibility:}
The primary threats to reliability and reproducibility are:
\begin{itemize}
    \item Non-deterministic outputs. Identical prompts and configurations can yield different outputs across runs due to factors such as floating-point arithmetic, batching, and stochastic decoding strategies.
    \item Infrastructure dependence. Results may vary depending on the hardware, software stack, and hosting environment used, making exact replication challenging across different infrastructure setups.
    \item Resource inequality preventing replication. LLM research is resource-intensive and hence remains predominantly in the domain of private companies or well-funded research institutions~\cite{schwartz2020green, ahmed2023industry}, excluding researchers from under-resourced institutions.
\end{itemize}

Researchers \must discuss the measures taken to increase reliability and reproducibility.
However, non-deterministic reproducibility is not inherently disqualifying. Qualitative traditions such as ethnography, grounded theory, and action research ensure trustworthiness through \emph{credibility}, \emph{transferability}, \emph{dependability}, and \emph{confirmability}~\cite{guba1981criteria}.
SaaS-based LLM research faces analogous conditions, as providers frequently deprecate model versions without guaranteeing stable behavior.

\paragraph{Ethical \& Regulatory Boundaries:}
The primary concerns for ethical and regulatory matters are:
\begin{itemize}
    \item Use of sensitive or proprietary data. Studies involving proprietary code, confidential business data, or personally identifiable information may face restrictions on data sharing that limit reproducibility.
    \item Jurisdictional obligations. Data protection regulations (e.g., GDPR, CCPA) and institutional policies may impose constraints on data collection, processing, and sharing in LLM-based studies.
    \item Implicit model bias. Especially for qualitative research LLM might \enq{reinforce dominant paradigms and biases} and \enq{identify, replicate and reinforce dominant language and patterns}~\cite{jowsey2025reject}.
\end{itemize}

Studies involving sensitive data \must discuss data governance mechanisms tailored towards LLM environments, compliant with applicable juristical obligations. If applicable, researchers \should discuss how model biases potentially impact the study outcomes and how those biases were evaluated.

\paragraph{Environmental \& Sustainability Constraints:}
\begin{itemize}
    \item Energy consumption. With growing model size, the environmental impact of experiments with LLMs increases. Considering the environmental impact in the study design is important given the substantial energy costs of LLM experiments~\cite{strubell2019energy}.
    \item Trade-off between repetition and sustainability. Increased reliability through repetition and the aim of cautious energy consumption pose an inherent tension field to navigate during the study design.
\end{itemize}

The study findings \should be discussed in the context of the resource-intensiveness of LLMs to justify the increased resource consumption over traditional approaches.

\paragraph{Mitigation Strategies}
The following mitigation strategies can help address the threats described above, depending on the study scope and feasibility.

\begin{itemize}
    \item \emph{Replication Packages} covering prompt and architecture specifications, model outputs, and representative examples for partial replicability, accompanied by an implementation using an open model for long-term stability.
    \item \emph{Human Validation} of subjective constructs, following quality criteria known from qualitative research.
    \item \emph{Longitudinal Re-Runs} to repeat experiments with LLMs over time, complemented by statistical analyses.
    \item \emph{Methodological Trustworthiness Measures.} Researchers \should consider triangulation, reflexivity, audit trails, and peer debriefing as complementary measures when deterministic reproduction is structurally impossible.
    \item \emph{Triangulation} via multiple models (e.g., proprietary and open models), multiple independent datasets, and multiple complementary metrics.
    \item \emph{Cost Accounting} through reporting of input/output tokens, token and service costs, or hardware specifications.
    \item \emph{Energy Preservation} by selecting smaller or newer less resource-intensive models and employing techniques such as input/output token reduction, model pruning, quantization, or knowledge distillation~\cite{mitu2024hidden} where feasible. Carbon footprint estimation is desirable, but still difficult.
    \item \emph{Ethical and Regulatory Considerations} through data governance mechanisms, ethical reviews, and bias assessment procedures.
    \item \emph{Sensitivity Analysis} through variation of LLM configurations, prompts, architecture decisions, datasets, and if applicable human participant backgrounds.
\end{itemize}

\guidelinesubsubsection{Study Types}

Researchers \must follow this guideline for all study types. Transparently reporting limitations and mitigations is a universal requirement, but specific concerns vary by study type.
For \annotators, researchers \must discuss potential biases in label assignment, label reliability limitations, and sensitivity of annotations to prompt wording and model choice.
For \judges, researchers \must address measurement validity concerns, known biases such as position bias or verbosity bias, and the extent to which LLM judgments align with human expert assessments.
For \synthesis, researchers \must discuss the risk of contextual misinterpretation, potential loss of nuance in summarized or aggregated outputs, and reflexivity limitations inherent in using an LLM for qualitative interpretation.
For \subjects, researchers \must discuss the fundamental inability of LLMs to truly simulate human behavior, the risk of stereotype amplification, and the limited ecological validity of simulated responses.
For \llmusage, researchers \must discuss generalizability constraints across different tools and user populations, and acknowledge how observed usage patterns may not transfer to other contexts.
For \newtools, researchers \must discuss replicability constraints arising from dependencies on commercial models, the impact of model updates on tool behavior, and limitations of the evaluation setup.
For \benchmarkingtasks, researchers \must discuss potential data contamination, benchmark scope limitations, and the extent to which benchmark performance generalizes to real-world tasks.

%% Suggestion: Consider softening "perfunctory or defensive" → "formulaic or insufficiently engaged with the study's specific threats".
%% Suggestion: Consider adding a bridging sentence connecting back to the "request additions first" principle used in other guidelines.
\guidelinesubsubsection{Advice for Reviewers}

Reviewers should verify that the limitation section is comprehensive and appropriate for the specific study type, checking that:
(1) limitations address the specific threats relevant to the study type (e.g., label reliability for annotation studies, simulation fidelity for studies using LLMs as subjects);
(2) mitigations are concrete and correspond to identified limitations rather than being generic statements;
(3) the impact of LLM non-determinism on findings is discussed;
(4) generalizability constraints---across models, configurations, time periods, and populations---are acknowledged.
When important limitations are missing, reviewers should request they be added. The absence of a limitation section, or one that is formulaic or insufficiently specific, is a more serious concern than any individual missing limitation and may warrant a major revision.

\guidelinesubsubsection{See Also}
\begin{itemize}[label=$\rightarrow$]
    \item Section~\modelversion: model configuration, cost accounting, and sensitivity analysis.
    \item Section~\toolarchitecture: architecture-level triangulation and sensitivity analysis.
    \item Section~\prompts: prompt reporting, audit trails, and longitudinal re-runs.
    \item Section~\benchmarksmetrics: metric selection and benchmark limitations.
    \item Section~\humanvalidation: human validation of subjective constructs.
\end{itemize} 
\begin{table}[tb]
\caption{Rationale and key recommendations for each guideline. \iconM~=~\must, \iconS~=~\should.}
\label{tab:rationale-recommendations}
\small
\begin{tabularx}{\textwidth}{@{}l>{\raggedright\arraybackslash}X>{\raggedright\arraybackslash}X@{}}
\toprule
\textbf{ID} & \textbf{Rationale} & \textbf{Core Recommendations} \\
\midrule
G1 & Transparency enables informed assessment of scope and limitations. &
\iconM Declare which LLM, how it was used, and where in the research process. \\
\addlinespace
G2 & Reproducibility requires precise identification of the system used in a study. &
\iconM Report exact version, date, configuration, and fine-tuning details. \\
\addlinespace
G3 & Complex systems require holistic documentation beyond the model itself. &
\iconM Describe full architecture, agent roles, and infrastructure. \\
\addlinespace
G4 & Prompts shape LLM behavior and must be verifiable. &
\iconM Publish all prompts, context files, and development strategy.
\newline \iconS Share interaction logs where feasible. \\
\addlinespace
G5 & Automated metrics alone cannot ensure validity of subjective constructs. &
\iconS Validate against human judgment with inter-rater reliability, where applicable. \\
\addlinespace
G6 & Reproducibility depends on access to the model under study. &
\iconS Include open LLM baseline; provide replication package. \\
\addlinespace
G7 & Meaningful evaluation requires reasoned valid measurement. &
\iconM Justify metric and benchmark choices; discuss their validity.
\newline \iconS Repeat experiments and report descriptive statistics. \\
\addlinespace
G8 & Honest acknowledgment of threats strengthens a study. &
\iconM Report threats by validity type.
\newline \iconS Employ mitigations where possible. \\
\bottomrule
\end{tabularx}
\end{table}
 
\section{Conclusion}

LLM non-determinism, opaque training data, and rapidly evolving models
threaten the reproducibility and replicability of empirical SE studies.
This paper presents a taxonomy of seven study types that organizes how
LLMs are used in SE research, together with eight guidelines for
designing and reporting such studies.
Each guideline distinguishes requirements (\must) from recommended
practices (\should) and is contextualized by the study types it
applies to.
Our guidelines recommend that researchers:
\begin{enumerate*}[label=(\arabic*)]
\item declare LLM usage and role;
\item report model versions, configurations, and customizations;
\item document the tool architecture beyond the model;
\item disclose prompts, their development, and, where feasible,
interaction logs;
\item validate LLM outputs with humans;
\item include an open LLM as a baseline;
\item use suitable baselines, benchmarks, and metrics; and
\item articulate limitations and mitigations, including costs, potential
biases, and environmental impact.
\end{enumerate*}
An applicability matrix maps guidelines to study types, and a reporting
checklist (Section~\ref{sec:checklist}) supports authors and reviewers.
The study types and guidelines were derived collaboratively by
22~co-authors, following a process similar to that of guidelines in
other disciplines (see, e.g.,~\cite{Begg1996}).

Following the quality dimensions of the \emph{SIGSOFT Empirical
Standards}~\cite{ralph2021empiricalstandardssoftwareengineering}, we
consider our guidelines
(1)~\emph{comprehensive}, as they cover seven study types that we
identified and were derived from 22 co-authors' collective expertise;
(2)~\emph{useful}, as they provide actionable recommendations with an
applicability matrix and a concrete reporting checklist
(Section~\ref{sec:checklist});
(3)~\emph{well-argued}, as each guideline is motivated by a dedicated
rationale and supported by references to the literature; and
(4)~\emph{integrated with previous work}, as they build on and
complement established empirical SE standards while addressing
LLM-specific challenges.
Adopting these practices will not eliminate LLM non-determinism, but it
should make claims auditable and results easier to replicate and compare
across studies.
Because models and tools evolve, we maintain the guidelines as a living
resource and invite the community to refine and extend them at
\href{https://llm-guidelines.org/}{llm-guidelines.org}.

\section*{Acknowledgments}

Our study types and guidelines are based on discussion sessions with researchers at the \emph{2024 International Software Engineering Research Network} (ISERN) meeting and at the \emph{2nd Copenhagen Symposium on Human-Centered Software Engineering AI} (supported by the Carlsberg Foundation with grant CF24-0693 and the Alfred P. Sloan Foundation with grant G-2024-22586 to Daniel Russo).
LLM-based tools, including ChatGPT (GPT-5) and Claude Code (Opus 4.6), were used interactively for language editing and structural revision suggestions across all sections of this paper. All LLM-generated suggestions were reviewed and revised by the authors.

\begin{appendices}

\section{Reporting Checklist}
\label{sec:checklist}

The following checklist, inspired by CONSORT~\citep{Schulz2010}, summarizes actionable items from the guidelines based on the \tldr sections.
The checklist is organized along typical paper sections.
Items marked \iconM are requirements (\must); items marked \iconS are recommendations (\should).
Each item references its source guideline (G1--G8).
Items annotated with \paper or \supplementarymaterial indicate where information should be reported; unmarked items may be reported in either.

\paragraph{\textbf{Introduction}}
\begin{itemize}[label=$\square$]
    \item \iconM Disclose any use of LLMs in the empirical study, specifying which LLM, how, and where it was used~(G1).
    \item \iconS Report the purpose of using LLMs, automated tasks, and expected benefits in the \paper~(G1).
\end{itemize}

\paragraph{\textbf{Research Design and Methods}}\mbox{}\\
\noindent\textit{Model Selection and Configuration}
\begin{itemize}[label=$\square$]
    \item \iconM Report the exact LLM model or tool version, configuration, and experiment date in the \paper~(G2).
    \item \iconM For fine-tuned models, describe the fine-tuning goal, dataset, and procedure in the \paper~(G2).
    \item \iconS Report default parameters and explain model and version choices~(G2).
    \item \iconS For quantized models, report the quantization level (e.g., 4-bit, 8-bit) and method (e.g., GPTQ or AWQ)~(G2).
    \item \iconS Compare base and fine-tuned models using suitable metrics and benchmarks; share fine-tuning data and weights as \supplementarymaterial (or justify in the \paper why they cannot be shared)~(G2).
\end{itemize}

\noindent\textit{Architecture}
\begin{itemize}[label=$\square$]
    \item \iconM Describe the full architecture of LLM-based tools in the \paper, including the role of the LLM, interactions with other components, and overall system behavior~(G3).
    \item \iconM For ensemble architectures, explain the coordination logic between models in the \paper~(G3).
    \item \iconM If autonomous agents are used, specify agent roles, reasoning frameworks, and communication flows~(G3).
    \item \iconM Report hosting, hardware setup, and latency implications~(G3).
    \item \iconM For tools using retrieval or augmentation methods, describe data sources, integration mechanisms, and update and versioning strategies~(G3).
    \item \iconS Include architectural diagrams and justify design decisions~(G3).
\end{itemize}

\noindent\textit{Prompts and Interactions}
\begin{itemize}[label=$\square$]
    \item \iconM Describe prompt development strategies (e.g., zero-shot, few-shot), rationale, and selection process in the \paper~(G4).
    \item \iconM Publish all prompts or, when using templates, prompt templates with representative instances, including their structure, content, formatting, and dynamic components, as \supplementarymaterial~(G4).
    \item \iconM Document input handling and token optimization strategies when prompts are long or complex~(G4).
    \item \iconM Report generation and collection processes for dynamically generated or user-authored prompts~(G4).
    \item \iconM Specify prompt reuse across models and configurations~(G4).
    \item \iconS If full prompt disclosure is not feasible, provide summaries or examples~(G4).
    \item \iconS Report prompt revisions and pilot testing insights~(G4).
    \item \iconS Include full interaction logs (prompts and responses) as \supplementarymaterial if privacy and confidentiality can be ensured~(G4).
    \item \iconS For agentic systems, report all context files (e.g., \texttt{AGENTS.md}) used to configure AI agents as \supplementarymaterial~(G4).
    \item \iconS For agentic systems, report developed plans and generalize interaction logs to include human-agent interaction traces as \supplementarymaterial~(G4).
\end{itemize}

\noindent\textit{Human Validation}
\begin{itemize}[label=$\square$]
    \item \iconM If using human validation, define the measured construct (e.g., usability, maintainability) and describe the measurement instrument in the \paper~(G5).
    \item \iconM When developing or adapting measurement instruments, share them~(G5).
    \item \iconS Consider human validation early in the study design and build on established reference models for human-LLM comparison~(G5).
    \item \iconS Validate LLM judgments against human judgment, report aggregation methods, and assess human-LLM agreement~(G5).
    \item \iconS Control for confounding factors and conduct power analysis to ensure statistical robustness~(G5).
\end{itemize}

\noindent\textit{Benchmarks and Metrics}
\begin{itemize}[label=$\square$]
    \item \iconM Justify all benchmark and metric choices in the \paper~(G7).
    \item \iconM Explain in the \paper why the selected metrics are suitable for the specific study~(G7).
    \item \iconS Include an open LLM as a baseline when using commercial models and report inter-model agreement~(G6).
    \item \iconS Summarize benchmark structure, task types, and limitations~(G7).
    \item \iconS Justify the number of experiment repetitions, for example through a power analysis or by monitoring convergence of descriptive statistics~(G7).
\end{itemize}

\noindent\textit{Reproducibility, Ethics, and Resources}
\begin{itemize}[label=$\square$]
    \item \iconM Provide raw LLM outputs for reproducibility and discuss sensitive data handling~(G8).
    \item \iconS Justify LLM usage in light of its resource demands~(G8).
    \item \iconS If full data sharing is not possible, include a subset of the validation data for partial replication~(G8).
    \item \iconS Provide a full replication package with step-by-step instructions as \supplementarymaterial~(G6).
\end{itemize}

\paragraph{\textbf{Results}}
\begin{itemize}[label=$\square$]
    \item \iconM If comparing models or tools, use appropriate inferential statistics (e.g., hypothesis tests, effect sizes) rather than relying solely on summary statistics~(G7).
    \item \iconS Repeat experiments due to the inherent non-determinism of LLMs and report the result distribution using descriptive statistics~(G7).
    \item \iconS Use traditional (non-LLM) baselines for comparison where possible~(G7).
    \item \iconS Report established metrics to make study results comparable; additional metrics may be reported where appropriate~(G7).
\end{itemize}

\paragraph{\textbf{Limitations and Threats to Validity}}
\begin{itemize}[label=$\square$]
    \item \iconM Describe measurement constructs and methods; disclose any data leakage risks and avoid leaking evaluation data into LLM improvement pipelines~(G8).
    \item \iconM Acknowledge non-disclosed confidential or proprietary components as reproducibility limitations~(G3).
    \item \iconM Transparently report study limitations, including the impact of non-determinism and generalizability constraints~(G8).
    \item \iconM Specify whether generalization across LLMs or across time was assessed, and discuss model and version differences~(G8).
    \item \iconS Employ and report strategies to mitigate identified validity and reproducibility threats, such as replication packages, human validation, longitudinal re-runs, triangulation, and sensitivity analysis~(G8).
\end{itemize} 
\end{appendices} 
\section{Statements and Declarations}

\paragraph{Competing Interests:} The authors declare no competing interests. Multiple co-authors are members of the EMSE Editorial Board, which is disclosed here for transparency.

%extra v-space will disappear when the "added" tag is deleted.
\paragraph{Authorship:} Each author contributed or reviewed content and, most importantly, read the complete guidelines and confirmed that they stand behind all recommendations, not only their own contributions.

\paragraph{Funding:} No funding was received for writing this article. However, the \emph{2nd Copenhagen Symposium on Human-Centered Software Engineering AI}, which played a central role in forming the team of co-authors, was supported by the Carlsberg Foundation with grant CF24-0693 and the Alfred P. Sloan Foundation with grant G-2024-22586 to Daniel Russo.

\paragraph{Ethical Approval:} No ethics approval was required for writing this article.

\paragraph{Data Availability:} This article has no associated data. The guidelines are, however, \href{https://llm-guidelines.org/}{hosted online} and \href{https://github.com/se-uhd/llm-guidelines-website}{maintained on GitHub}.

\bibliographystyle{spbasic}
\bibliography{literature}

\end{document}
